% UNINTEGRATED ARITHMETIC15 FOUNDATION PART, 25 October2026.
% Same existing arithmetic paper family; source preparation only.
% First batch: accepted general dihedral3/19 theorem and D14 consumer.
% DRAFT pending independent and coordinating review; no execution/compilation.
\section{A strengthened ordinary dihedral comparison}
\label{app:arithdihedral}

This section supplies the ordinary rank-two comparison needed beside
\(D_{14}\).  Its component theorem is general real geometry.
Arithmeticity enters only when the already established arithmetic
\(D_{14}\) prefix bound is used in the final corollary.
Throughout, \(p\) is the ORIGINAL vertex farthest from one strict
interior point.  The theorem is not asserted at an arbitrary
neighbour with a weakened comparison factor.

\begin{proposition}\label{prop:ordinarydihedralthreenineteenths}
Let a compact hyperbolic Coxeter polytope have an actual farthest
vertex \(p\), and write \(z=E(p-x)\), \(\|x\|<1\).
Let one WHOLE component of the incident finite diagram be
\(I_2(m)\), \(m\ge4\), and let \(e\) be the squared norm of the
projection of this SAME \(x\) onto that component.
If \(e\le3/19\), then \(m=7\) and one actual opposite wall
closes both component directions, has a single old contact of
label three there, and has no contact with any other old component.
This is precisely the actual \(O_7\) profile.
Consequently every ordinary \(I_2(m)\), \(m\ge4\), has energy
strictly greater than \(3/19\), with arbitrary finite spectators.
\end{proposition}

\begin{proof}
Use outward unit incident roots, finite Gram \(B\), inverse
\(C=B^{-1}\ge0\), duals \(\omega_i\), and unit inward dual
directions \(-u_i=-\omega_i/\sqrt{C_{ii}}\).
For an actual outgoing edge of length \(L_i\), maximality at
the original \(p\) gives
\[
 x\cdot u_i\ge\tau_i=\tanh(L_i/2)>0,\qquad
 r_i=\tanh L_i={2\tau_i\over1+\tau_i^2}.
\]
The old dihedral inverse is
\[
 C_I=d\begin{pmatrix}1&k\\k&1\end{pmatrix},\qquad
 k=\cos(\pi/m)\ge1/\sqrt2,\quad d=(1-k^2)^{-1}.
\]
Thus each \(\tau_i^2\le e\), and the increasing function
\(4t/(1+t)^2\) on \(0\le t<1\) gives
\[
                         r_i^2\le57/121=:r .
\tag{Dih1}\label{eq:arithdihedralr}
\]

We recall the actual-wall calculation forcing a common opposite.
For an outside actual wall \(c=Ap-v\), let
\(\beta_j=-c\cdot n_j\ge0\) and put
\[
 T_i(c)={(C\beta)_i\over A\sqrt{C_{ii}}}.
\]
For the actual opposite \(c_i\) of edge \(i\),
\(T_i(c_i)=1/r_i\).
The matrix
\(C-C_{\cdot j}C_{j\cdot}/C_{jj}\)
is entrywise nonnegative: its nonzero block is the inverse of
the positive principal Gram obtained by deleting \(j\).
Actual acuteness between two distinct opposites therefore gives
\(T_j(c_i)\le r_j\), while that same nonnegative matrix gives
\(T_j(c_i)\ge k/r_i\) within this old component.
Hence distinct opposites imply \(r_ir_j\ge k\).
But~\eqref{eq:arithdihedralr} gives
\(r_ir_j\le57/121<1/\sqrt2\le k\).
One actual wall \(c\) therefore closes BOTH directions.

Write its full decomposition as
\[
 c=Ap-v_I-v_S,\quad h=A^2>0,\quad
 b=-c\cdot n_1\ge0,\quad g=-c\cdot n_2\ge0,\quad
 Q=\|v_S\|^2\ge0 .
\]
Here \(S\) is the entire orthogonal spectator forest.
Every positive contact is a native magnitude at least \(1/2\);
none of its finite, parallel or dotted possibilities is omitted
before the actual endpoint constraints are imposed.
At least one of \(b,g\) is positive.

Suppose both are positive.
At the actual endpoint releasing direction one, the retained
white Gram has residual \(\delta=1-Q-g^2>0\).
Its actual edge formula is
\[
 \cosh L_1={\sqrt d(b+kg)\over\sqrt{1-Q-g^2}}
 \ge{\sqrt d(1+k)\over\sqrt3}
 ={\,\cot(\pi/(2m))\over\sqrt3}
 \ge{1+\sqrt2\over\sqrt3}>{11\over8}.
\]
The last comparison is \(128\sqrt2>171\), whose positive
squared margin is \(32768-29241\).
It follows that
\[
 e\ge\tanh^2(L_1/2)>
 {11/8-1\over11/8+1}={3\over19},
\]
a contradiction.  Hence exactly one old contact \(\beta>0\)
remains.

We need one whole-spectator fact at small time.
If \(h<1/4\), then \(c\) cannot contact any foreign component.
Indeed a positive foreign contact gives a normalized derivative
\[
 {(C_S\beta)_j\over\sqrt{(C_S)_{jj}}}\ge1/2 .
\]
This uses the full spectator inverse, with nonnegative entries
and diagonal at least one.
Its ratio to \(\sqrt h\) exceeds one.
The actual-wall shadow inequality therefore forces this SAME
\(c\) to be the opposite of that foreign edge as well.

For clarity, the required orthogonal obstruction is specific to
the original farthest vertex.  Suppose a wall closed orthogonal
dual directions \(u_i,u_j\).
Clear it positively off all other incident roots \(W\):
\[
 c_F=c+\sum_{s\in W}\alpha_sn_s,\qquad
 \alpha=B_W^{-1}(-c\cdot n_W)\ge0.
\]
Strict interiority gives \(z\cdot c_F<0\).
In the retained Lorentz three-space its exact form is
\[
 c_F=A p-(A/r_i)u_i-(A/r_j)u_j .
\]
Therefore
\[
 {x\cdot u_i\over r_i}+{x\cdot u_j\over r_j}<1.
\]
Original maximality gives the opposite bound
\(\tau_i/r_i+\tau_j/r_j
 =1+(\tau_i^2+\tau_j^2)/2>1\).
This proves the obstruction, including arbitrary contacts in \(W\).
Thus \(h<1/4\) really forces \(Q=0\).

Put \(q=d\beta^2\).
The smaller normalized derivative square is \(k^2q\).
The common-wall edge formula and~\eqref{eq:arithdihedralr} imply
\[
 q-1\le h=q+Q-1\le rk^2q .
\]
Substituting \(k^2=1-1/d\) and \(r=57/121\) gives the exact
strengthened restriction
\[
                         \beta^2(64d+57)\le121 .
\tag{Dih2}\label{eq:arithdihedralcut}
\]
This is derived at the \(3/19\) threshold; it is not the older
\(1/7\) inequality with an unchanged coefficient.

Since \(\beta^2\ge1/4\), this forces \(d\le427/64\).
For every old \(m\ge8\),
\(d\ge4+2\sqrt2>34/5>427/64\), impossible.
The complete remaining old alphabet is \(4,5,6,7\).
At old \(m=5\),
\(d=2+2/\sqrt5>185/64\):
equivalently \(128>57\sqrt5\), with squared margin139.
Since \(d\) increases with \(m\), if
\(\beta\ge1/\sqrt2\) at any of old \(5,6,7\),
\eqref{eq:arithdihedralcut} would give the opposite
\(d\le185/64\).
Thus all three rows have exactly \(\beta=1/2\).
Every larger finite contact, parallel and dotted contact has
been excluded by the native jump, rather than a sampled label list.

At old four, \(d=2\) and
\(\beta^2\le121/185<\cos^2(\pi/5)\).
The strict comparison is \(413<185\sqrt5\), with squared
margin556.  Only contacts three and four remain, giving
\(q=1/2\) or \(1\).
Consequently \(0<h\le57/242<1/4\).
The small-time argument just proved gives \(Q=0\), whereas
\(h=q-1\le0\), a contradiction.

At old five only \(\beta=1/2\) remains.
Since \(d<3\), \(q<3/4\) and
\(0<h\le r(d-1)/4<57/242<1/4\).
Again \(Q=0\), and \(h=q-1<0\).

At old six, \(\beta=1/2\), \(d=4\), \(q=1\).
If \(Q=0\), then \(h=0\), impossible.
For any \(Q>0\), the actual positive sum
\[
                          \nu=c+2n_1+\sqrt3n_2
\]
is a nonzero future null: all three selected products vanish.
The coefficient order is for OLD6/NEW3, since
\(c\cdot n_1=-1/2\), \(n_1\cdot n_2=-\sqrt3/2\),
\(c\cdot n_2=0\).
Every other actual wall pairs nonpositively termwise with \(\nu\).
It is a feasible ideal ray, contrary to compactness.
No foreign contact was discarded in this affine exclusion.

At old seven only \(\beta=1/2\) remains, and
\[
 0\le Q=1+h-d/4
 \le{427-64d\over484}<{107\over484}<{1\over4},
\]
using \(d_7>5\).
Every positive spectator cost is at least \(1/4\), by inverse
positivity and the native contact lower bound.  Thus \(Q=0\).
We have exactly the ACTUAL isolated-new-three wall in old
\(I_2(7)\), closing both directions: the stated \(O_7\).
Its time square is \(d_7/4-1\in(0,1/3)\);
the bounds \(5<d_7<16/3\) follow from
\(6/5<2\cos(2\pi/7)<5/4\).
Its separate inherited energy floor \(>9/100\) is unchanged.

All weak inequalities included \(e=3/19\), so no ordinary equality
case is omitted.  The proof has retained every finite spectator
component, all its contacts and the same original \(x\).
No arithmetic hypothesis or weakened-neighbour comparison was used.
\end{proof}

\begin{corollary}\label{cor:D14ordinarydihedral}
At the original farthest vertex in arithmetic dimension sixteen,
a whole \(D_{14}\) cannot have an ordinary irreducible
rank-two spectator.
\end{corollary}
\begin{proof}
The fresh true-rank-\(14\) argument in the proof of
Proposition~\ref{prop:f16} gives
\[
                         e_D>.84816={5301\over6250}
\]
for its actual whole rank-two spectator.
Its bootstrap uses only the independently available spectator
floor \(>.09\), so it does not depend on the present strengthening.
For an ordinary \(I_2(m)\), \(m\ge4\), the proposition gives
\[
 \|x\|^2=e_D+e_S>
 {5301\over6250}+{3\over19}
 =1+{719\over118750}>1 .
\]
For ordinary \(A_2\), use its geometric companion bound
\(319/2000\), which exceeds \(3/19\) by \(61/38000\).
Both contradict \(\|x\|<1\).
The actual typed \(N_7/O_7\) alternatives retain their separate
costs and arguments; they have not been charged as ordinary.
\end{proof}

% Next batches may append reviewed arithmetic proof bodies here.
% Preserve this exact first-batch source before any such expansion.

\section{Arithmetic chain windows and simultaneous actual feet}
\label{app:arithmeticchain}

The endpoint arguments below share one arithmetic chain calculation.
We prove it before specializing to the whole \(A_{16}\) and
\(A_{17}\) cases.  The geometric finite-white classification is the
one in the general companion; the arithmetic restrictions and the
simultaneous treatment of isolated feet are given here.
No component dimension is obtained by deleting a node from a
previous ambient theorem.

Let \(N=16,17,18\), let the original incident forest contain the
WHOLE component \(A_{N-1}\), and retain all finite spectators.
Use its zero-sum \(\mathbb R^N\) frame and unit old roots
\(n_i=(e_i-e_{i+1})/\sqrt2\).  Then
\[
 D_i=\omega_i^2={2i(N-i)\over N},\qquad
 C_{ij}=\omega_i\cdot\omega_j={2i(N-j)\over N}\quad(i\le j).
\]
All heights \(h_i=x_A\cdot\omega_i\) concern the SAME original
farthest vector.  Set
\[
             \ell_{16}=.601,\qquad\ell_{17}=\ell_{18}=.6.
\]

\subsection{The complete central inventory at the new losses}
\label{app:arithchaininventory}

For \(5\le i\le N-5\), call an index weak when
\(h_i\le\sqrt{D_i}-\ell_N\).
We claim that every central constraint is either strong, an
isolated-within-\(A\) cap with its actual cleared parameter
\(0<u<1\), or a pure golden end-helper.
The last alternative can occur at at most ONE central index.

First establish the full inventory, without assuming a contrary
component energy.  Deleting \(i\) retains two old paths of ranks
at least four.  Their complete one-node finite extensions are
isolation; one admissible end, \(D\)- or \(E\)-position three;
one terminal four; or two endpoint threes, one in each path.
Two contacts in a single path make a cycle.
A connected two-path completion has rank \(N-1>8\), so finite
classification forces the ordinary path bridge, with no extra
finite foreign arm.

The zero-own repair depends only on side order \(m\ge5\).
At a far-end three the time square is less than its full positive
foreign cost; the next-to-far three has foreign remainder \(<2/m\)
and the same obstruction.
If such a positive foreign contact existed, its actual shadow
would force the same opposite to close an orthogonal component,
contrary to the original-farthest obstruction proved above.
Terminal four has foreign remainder \(<1/m\le1/5\), so no positive
spectator cost, and then negative time.
The other zero-own \(D/E\), internal-four and two-white roles
have their actual positive affine stars or cycles.
Their time is positive through the actual cap and every omitted
wall product is nonpositive.  Own contact is therefore positive.

Write \(q=\sqrt{D_i}\).  Here \(5/2<q\le3\).
The actual edge ratio needed for a loss \(\ell\) is
\[
 R_\ell(q)={1\over\ell}
                -{q-\ell\over q(2q-\ell)}.
\]
It increases with \(q\) and decreases with \(\ell\) on this
rectangle.  Thus
\[
 R_\ell(q)\le R_{.6}(3)=41/27<761/500.
\]
The finite-white residual formulas give the following ordinary
lower ratios, with side order at least five:
near-three/own at least five has
\((101/125+2/5)/\sqrt{3/5}\);
a \(D\)-position has
\((101/125+1/5)/\sqrt{2/5}\);
the \(E\), terminal-four and two-end-three rows are stronger.
All exceed \(761/500\).
The smaller owns have the same actual affine principals.
For the two exceptional roles WITH positive foreign cost,
the complete clearings give ratio
\(>227/(25\sqrt{35})>761/500\), or ratio square \(>121/35\).
These estimates retain the ENTIRE contacted spectator cost.
Hence no additional ordinary nonisolated row is weak.

Unlike the older \(3/2\) test, the present ratio threshold retains
near-own-four also at side order seven.  Up to reflection all
remaining such profiles are
\[
          v=\omega_i/\sqrt2+\omega_{i-1}/2,\qquad i=5,6,7 .
\]
The simple-three helper normalizes the cap to square two over the
actual field, with old product \(-1\); its own product is
\(-\sqrt2\), so \(\sqrt2\in K\).
At the embedding reversing \(\sqrt2\), the time numerators at
\(N=16\) are
\[
        126-88\sqrt2,\quad143-100\sqrt2,\quad154-108\sqrt2 .
\]
They are positive since \(\sqrt2<17/12\).
Their increments with \(N\) are
\((i-1)(3-2\sqrt2)>0\).
The conjugate new-root Schur complement is the negative of this
numerator divided by \(N\), contradicting positive definiteness.
This removes all three side orders and both orientations.

The remaining nonisolated possibility is pure:
\[
 c=Ap-\gamma\omega_i-\omega_e/2,\qquad
 e=1\text{ or }N-1,\quad Q=0,
\]
with own five, finite own at least seven, or dotted.
The exact helper-side identity is
\[
 {(\gamma+1/(2m))^2\over(m+1)/(2m)}-(4\gamma-2)
 ={(\gamma-1-1/(2m))^2\over(m+1)/(2m)}\ge0.
\]
The ratio threshold and native-five gap therefore give
\[
             202/125<t=2\gamma<3139/1458<54/25.
\tag{Astrip}\label{eq:arithAhelperphysical}
\]

For isolated-within-\(A\) contact, retain the whole foreign
projection \(v_S\) and write
\[
 c=Ap-\gamma\omega_i-v_S,\quad Q=\|v_S\|^2,\quad
 u=(1-Q)/\gamma^2,\quad A^2=\gamma^2(D_i-u).
\]
Here \(Q<1\) by finite-white positivity.
When \(Q>0,u>1\), the actual positive clearing has loss at most
\[
 {1\over\sqrt{D_i}+\sqrt{D_i-1}}+\sqrt{u-1}
 <.2+.383<\ell_N .
\]
This uses \(Q\ge1/4\), \(\gamma\ge\varphi/2\), and
\(u-1\le5-3\varphi<.383^2\).
At \(u=1\) the positive cleared vector is a forbidden feasible
future null.
The interval \(0<u<1\) is RETAINED, including \(Q>0\);
we do not charge a strong floor there.

For pure isolated owns, three, four, six and parallel have their
actual affine/null principals.
Own five has the actual \(H_4\) projected-point loss
\(<25/76+5/24<.6\).
For own seven, the increasing cubic
\(8a^3-4a^2-4a+1\) gives \(.9<a<.902\).
The actual positive \(A_2\) core gives
\[
 h_i>\sqrt{D_i-a^{-2}}-\sqrt{4/3-a^{-2}},
\]
and hence loss
\[
 <{100/81\over2.6+2.37}+\sqrt{17/150}
 <.25+.337=.587 .
\]
Pure own eight is arithmetically impossible:
\(4\gamma^2=2+\sqrt2\), and its opposite conjugate would require
\(D_i(2-\sqrt2)<4\), whereas
\(D_i\ge55/8>4+2\sqrt2\).
Every embedding of this minimal field extends to the actual
totally real field, so this contradiction applies there.
For all finite owns at least nine, the exact singleton finite
projection gives
\(h_i>\sqrt{D_i-1}-\tan(\pi/m)\).
Since \(\tan(\pi/9)<1386/3727<2/5\), the loss is \(<.6\).
This covers the infinite finite-own tail, not a finite label sample.
All positive finite and negative-norm projections just used are
positive sums of ACTUAL roots; their selected signs, positive
time and all omitted walls are retained as in the companion's
actual-point lemmas.

\subsection{Every pure weak helper is golden over the actual field}
\label{app:arithgoldenfield}

Normalize old \(A\) roots to square two.
If \(r\) is a \(K\)-root on the actual cap line and its simple-three
helper is \(a_e\) of square two, then
\(r^2=2(r,a_e)^2\).  Rescaling by \(-(r,a_e)^{-1}\in K\)
therefore makes the cap square two with helper product \(-1\).
Its own product is \(-t\).
The path cycle through old simple-three bonds makes \(t\) integral.
This compares only connected root lines; no disconnected norm
squareclasses are equated.

Every nonidentity embedding of \(\mathbb Q(t)\) extends to \(K\).
Strict conjugate positivity of the full old-component Schur
complement gives
\[
                    D_it_\sigma^2+2C_{ie}t_\sigma+D_e<4 .
\]
After reversing \(i\le N/2\), the cheaper costs for negative and
positive \(t_\sigma\) are respectively
\[
 Q_-(s)=2is^2+4s+2-2(is+1)^2/N,\qquad
 Q_+(s)=2is^2+2-2(is-1)^2/N .
\]
At the smallest \(N=16,i=5\), their values at \(-4/5,12/25\)
exceed four by \(3/40,59/1000\), respectively.
The costs at those endpoints increase with \(i\) and \(N\);
their values at zero are below four.
Convexity of the quadratics therefore gives the full strict strip
\[
                         -4/5<t_\sigma<12/25 .
\tag{Agoldenstrip}\label{eq:arithAgoldenstrip}
\]
This covers EVERY nonphysical conjugate, not chosen ones.

If the degree is at least five, the nonzero integral norm has
magnitude \(<(54/25)(4/5)^4<1\).
At degree four or three its magnitude is below two.
Any positive small conjugate would make the magnitude below one,
so all small conjugates are negative and their absolute product
is \(1/t\).
At degree four the integral trace lies in \((-1,0)\):
its lower bound is \(t-12/5>-1\), and its upper bound
\(t-3t^{-1/3}<0\) follows from \((54/25)^4<27\).
At degree three the trace lies in \((0,1)\), since
\(t-8/5>0\) and \(t-2/\sqrt t<1\);
the latter follows from
\((54/25)(29/25)^2<4\).
Both degrees are impossible.
At degree two the integral trace is one or two.
Trace two makes the nonzero norm \(t(2-t)\) lie in \((-1,1)\).
Trace one forces \(t<9/5\) and \(t(1-t)\in(-2,0)\);
the norm is \(-1\), and \(t=\varphi\).
Degree one leaves only \(t=2\), the actual parallel pair with
a feasible future null.
Thus every pure weak helper has \(\gamma=\varphi/2\).

The old component and this foreign-free cap have a physically
Lorentzian Gram over \(\mathbb Q(\varphi)\).
If \(K\) were a proper extension, a nonphysical real embedding
fixing that subfield would preserve this indefinite principal
Gram, impossible.
Consequently the helper forces the ACTUAL field
\(K=\mathbb Q(\varphi)\), rather than a convenient overfield.

\subsection{The complete golden packing calculation}
\label{app:arithgoldenpacking}

For a golden cap reversed to \(5\le i\le N/2\), its time
numerator \(M=2NA^2\) is
\[
 M_i^{\rm near}=i(N-i)-(N+1)+(i+2)(N-i)\varphi,\quad
 M_i^{\rm far}=i(N-i)-(N+1)+i(N-i+2)\varphi .
\]
These formulas retain BOTH helper ends.
At \(N=18\), conjugate positivity excludes the midpoint nine
and opposite-end indices six through nine.
The endpoint costs are \((173-84\varphi)/36\) and
\((197-99\varphi)/36\), both greater than one, and the
opposite-end cost is concave in \(i\).
At \(N=17\), only far seven is excluded, by
\((170-84\varphi)/34>1\) using \(\varphi<34/21\).
At \(N=16\) retain every orientation.
The resulting COMPLETE necessary tables are
\[
\begin{array}{c|c|r}
 N&\text{type and }M&\mathrm N(M)\\ \hline
16&\mathrm{near}\ 5:\ 38+77\varphi&-1559\\
 &\mathrm{near}\ 6:\ 43+80\varphi&-1111\\
 &\mathrm{near}\ 7:\ 46+81\varphi&-719\\
 &\mathrm{near}\ 8:\ 47+80\varphi&-431\\
 &\mathrm{far}\ 5:\ 38+65\varphi&-311\\
 &\mathrm{far}\ 6:\ 43+72\varphi&-239\\
 &\mathrm{far}\ 7:\ 46+77\varphi&-271\\ \hline
17&\mathrm{near}\ 5:\ 42+84\varphi&-1764\\
 &\mathrm{near}\ 6:\ 48+88\varphi&-1216\\
 &\mathrm{near}\ 7:\ 52+90\varphi&-716\\
 &\mathrm{near}\ 8:\ 54+90\varphi&-324\\
 &\mathrm{far}\ 5:\ 42+70\varphi&-196\\
 &\mathrm{far}\ 6:\ 48+78\varphi&-36\\
 &\mathrm{far}\ 8:\ 54+88\varphi&-76\\ \hline
18&\mathrm{near}\ 5:\ 46+91\varphi&-1979\\
 &\mathrm{near}\ 6:\ 53+96\varphi&-1319\\
 &\mathrm{near}\ 7:\ 58+99\varphi&-695\\
 &\mathrm{near}\ 8:\ 61+100\varphi&-179\\
 &\mathrm{far}\ 5:\ 46+75\varphi&-59 .
\end{array}
\]
At \(N=16\), near and far eight are reflections at the same
central index, so one row suffices.

Why do these finite tables restrict coexistence?
Subtracting a cap's \(K\)-valued projection onto the FULL old
forest leaves a nonzero vector on the SAME one-dimensional
time line over \(K\).
The caps have no foreign component.
Thus any two time coefficients have ratio in \(K\);
their \(M\)-ratio must be a square IN \(K\), and its rational
norm must be a rational square.
This argument retains all spectators.

At \(N=16\) and \(18\), the positive norm magnitudes in each
table are pairwise coprime nonsquares; their ratios are therefore
not rational squares.  These are direct small integer checks of
the displayed complete tables.
At \(N=17\), different norm squareclasses \(-1,-19,-179\)
cannot mix.
In class \(-1\), write the four entries as
\(42\varphi^3,18\varphi^5,14\varphi^5,6\varphi^7\).
After squares are removed the rational coefficient classes are
\(42,2,14,6\).
Every pair ratio has an odd valuation at three or seven;
none is a rational square or five times a rational square.
Those are exactly the positive rational squares in
\(\mathbb Q(\sqrt5)\), as expansion of
\((a+b\sqrt5)^2\) shows.
For the repeated \(-19\) class,
\[
 {M_{\rm far8}\over M_{\rm near6}}
   ={25+33\varphi\over76}={83+33\sqrt5\over152}
\]
has norm \(1/16\).
A square would require
\(a^2=121/304\) or \(45/304\), both with odd denominator
valuation at nineteen.
Thus different types cannot coexist, including the repeated
norm classes where a norm-only test would be insufficient.

At a fixed index the actual opposite is unique.
Reflected copies of the same noncentral type would have
\[
 -2c_L\cdot c_R=
 { (N-2i)(i+(i+2)\varphi)-(N+2)\over N}
\]
for near helpers, and the same formula with \(i-2\) replacing
\(i+2\) for far helpers.
Their simple-three helper bonds and old path make this
parameter integral by the invariant cyclic criterion, even for
a dotted mutual.
But the constant numerators for \(i=5,6,7\) at \(N=16\) are
\(12,6,-4\); for \(i=5,6,7,8\) at \(17\) they are
\(16,11,2,-11\); and at \(18\) they are \(20,16,8,-4\).
The far rows have the same respective constants.
None is divisible by \(N\), contrary to
\(\mathcal O_{\mathbb Q(\varphi)}=\mathbb Z[\varphi]\).
The central \(N=16\) row cannot occupy two different indices.
Hence at most ONE central index has a golden helper.
This is an all-native-five packing statement, independent of
which of those caps are actually weak at the new loss.

\subsection{Sum every retained isolated foot in one reference}
\label{app:arithsummedfeet}

The intervals \(0<u<1\) retained earlier cannot be declared
strong by the older \(.61\) estimate.
For such a cap clear the entire spectator projection by its
positive actual-root combination.  The actual vector
\[
 \eta=c+v_S+\gamma u n_i
      =\gamma\{\sqrt{D_i-u}\,p-(\omega_i-un_i)\}
\]
has square \(-\gamma^2u(1-u)\), positive time, zero contacted
spectator products, negative or zero \(c,n_i\) products, and
all other actual products nonpositive.
Its normalization is a genuine point of the whole chamber.
Original maximality gives
\[
 x_A\cdot(\omega_i-un_i)
 \ge\sqrt{D_i-u}-\sqrt{u(1-u)}.
\]
If \(\lambda_i=\sqrt{D_i}-\sqrt{D_i-u}+\sqrt{u(1-u)}\)
denotes its loss, then
\[
 \lambda_i\le a_i u+\sqrt{u(1-u)},\qquad
 a_i={1\over\sqrt{D_i}+\sqrt{D_i-1}}.
\]
The rational constants needed below are
\[
\begin{array}{c|ccc}
N&a_i<a&\text{positive reference coefficient }\beta_i\ge b&
 \ell\\ \hline
16&.1982&.022&.601\\
17&.1955&.0205&.6\\
18&.194&.018&.6 .
\end{array}
\]
The \(a\) bounds are positive-square comparisons at
\(D_{\min}=55/8,120/17,65/9\).
For every \(0<u<1\),
\[
 a u+\sqrt{u(1-u)}-\frac b2(2u-u^2)<\ell.
\tag{Aabsorb}\label{eq:arithAabsorb}
\]
For \(u\le1/2\) use \(a/2+1/2\), giving
\(.5991,.59775,.597\).
For \(u\ge1/2\), use \(2u-u^2\ge3/4\) and
\(\max(au+\sqrt{u(1-u)})=(a+\sqrt{1+a^2})/2\).
Radical upper bounds \(1.0195,1.019,1.019\) give respectively
\(.6006,.5995625,.59975\), all strictly below their \(\ell\).

Let \(H\) be a concave reference with available nonnegative
coefficients \(\beta_i\), positive ones at least \(b\):
\[
 w_j=(H_j-H_{j-1})/\sqrt2,\quad
 \beta_i=w\cdot n_i,\quad
 w=\sum_i\beta_i\omega_i,\quad
 E_H=\|w\|^2=\sum_i\beta_iH_i .
\]
Require \(H_i+\ell<\sqrt{D_i}\) at its positive central
coefficients.  Omit the possible golden index.
If \(S\) is the set of retained isolated indices with
\(\beta_i>0\), define ONE vector
\[
                       Z=w-\sum_{i\in S}\beta_i u_i n_i .
\]
Sum all actual point pressures BEFORE estimating its norm:
\[
 x_A\cdot Z\ge E_H-\sum_{i\in S}\beta_i(\lambda_i-\ell).
\]
Because \(w\cdot n_i=\beta_i\) and
\(n_i\cdot n_j\le0\) for \(i\ne j\),
\[
 \|Z\|^2\le E_H-\sum_{i\in S}\beta_i^2(2u_i-u_i^2).
\]
It follows that
\[
 e_A\ge2x_A\cdot Z-\|Z\|^2
 \ge E_H+\sum_{i\in S}\beta_i
       \{\beta_i(2u_i-u_i^2)-2(\lambda_i-\ell)\}>E_H.
\tag{Asummed}\label{eq:arithAsummed}
\]
Strictness follows from~\eqref{eq:arithAabsorb} if \(S\ne\varnothing\),
and from the strict available constraints otherwise.
Adjacent isolated indices and positive foreign costs are allowed.
No nonadjacency is assumed for this new isolated set, and no
independently optimized energies are added.

The references used by the endpoint arguments are
\[
\begin{array}{c|c|c|c}
N&H_5,H_6,\ldots\text{ to midpoint}&E_H&
 \text{maximum internal golden omission loss}\\ \hline
16&2.021,2.137,2.205,2.227&.8354522&.000576\\
17&2.056,2.186,2.269,2.310&.8708972&.00055225\\
18&2.087,2.228,2.324,2.381,2.399&.9037838&.00050625 .
\end{array}
\]
Reflect \(H_{N-i}=H_i\), use linear end tails, and at \(N=17\)
take \(H_9=H_8\).
The positive curvatures have exactly the lower bounds \(b\) above.
All floor inequalities are positive-square comparisons:
the new heights are the old \(.61\) heights shifted by
\(.61-\ell\), so \(H_i+\ell\) is unchanged.
A single internal omission sets its coefficient zero by neighbour
averaging and increases the adjacent coefficients.
An end omission replaces the first tail by the chord \(0\)--\(6\),
or its reflection.  These are changes in ONE reference.
In particular the full and omitted references preserve the
summed-foot argument whenever their positive central coefficients
have the displayed lower bounds; subsequent endpoint grafts must
check that condition again rather than assume it.

\subsection{A golden end hole pays for its own endpoint}
\label{app:arithgoldenend}

For later proper-component consumers we record how the exceptional
end omission is closed.  Suppose its central index is five,
reversing if necessary, and put \(\gamma=\varphi/2<.81\).
For either helper end \(e\), with its adjacent old node \(j\),
the actual finite \(H_4\) correction is
\[
 b_c=\gamma n_5+\tfrac23n_e+\tfrac13n_j,\qquad
 \eta=c+b_c,\quad
 \eta^2=\tfrac23-\gamma^2<.111^2 .
\]
The actual positive finite projected point has pure spatial
displacement \(v-b_c\).
Its normalization and original maximality give
\(x\cdot(v-b_c)\ge A-.111\).
Since \(x\cdot b_c>0\), the helper-end signed coordinate satisfies
\[
                  x_{\rm end}>
             \sqrt2(A-.111-\gamma h_5)>.49 .
\]
Both orientations are covered: the smaller opposite-end time
squares, at \(N=16,17,18\), are
\[
 (38+65\varphi)/32,\quad(42+70\varphi)/34,\quad
 (46+75\varphi)/36.
\]
They give \(A>2.115,2.136,2.155\).
Weakness gives \(h_5<2.022,2.057,2.088\).
Using \(\sqrt2>1.4\), the corresponding signed-coordinate lower
bounds are \(.512652,.502362,.493808\).
The nearer-end time is larger.
Here the signed coordinate is \(x_1\) at helper one and
\(-x_N\) at helper \(N-1\).
No strong floor at the missing index five has been inferred.

Use the chord \(0\)--\(6\) at the omitted end.
At the helper end the resulting reference is flat on \(L=6\)
coordinates in the near case and \(L=5\) in the far case,
with signed mean \(.25<b_0<.3\).
Replace that block by
\[
             .49,\quad
             {Lb_0-.49\over L-1},\ldots,
             {Lb_0-.49\over L-1}.
\]
Its sum is unchanged.  The latter coordinates lie in
\((.19,.262)\); the next old coordinate is below \(.101\).
Thus concavity persists, the new central boundary coefficient
is \(>.062>b\), and all other positive central coefficients
keep their previous lower bound.
The new endpoint coefficient is positive and is paid by the
actual endpoint pressure just proved.
The one-reference energy increases by
\[
                {L\over L-1}(.49-b_0)^2>.04332 .
\]
The pre-graft chord energy is exactly
\[
 E_H-\frac12\{H_5^2/5+(H_6-H_5)^2-H_6^2/6\}.
\]
At \(N=17,18\) it is respectively
\(12569249/15000000\) and \(13079276/15000000\);
at \(N=16\) it exceeds \(.8008\).
These direct rational formulas give, after the graft and
simultaneous absorption~\eqref{eq:arithAsummed}, the conservative
component bounds
\[
                 e_{A_{15}}>.8348,\qquad
                 e_{A_{16}}>.8703,\qquad
                 e_{A_{17}}>.9032 .
\]
If the single golden hole is internal instead, subtract the
internal losses in the reference table:
the energies are \(>.8348762,.87034495,.90327755\).
If there is no hole they are stronger.
The three conclusions therefore include every omission location,
all retained isolated feet, and arbitrary finite spectators.
They are the shared component estimates; no historical
ambient-dimension consumer has been imported.

\section{The actual whole-\(A_{17}\) endpoint reduction}
\label{app:wholearithA17}

We now prove the geometric and arithmetic inclusion needed before
using the recorded finite \(A_{17}\) propositions.
The old incident diagram is WHOLE \(A_{17}\), so \(N=18\) and
there is no spectator.
Keep the original \(p,z,x\) throughout.
Use the \(N=18\) reference and simultaneous-foot argument of
Appendix~\ref{app:arithmeticchain}, with energy
\(E_0=.9037838\), central coefficients at least \(.018\),
and first/last five-coordinate signed means
\(.2951<b<.2952\).
The next signed coordinate is below \(.1\).
A possible single golden central hole is omitted before its
pressure is estimated.

\subsection{Actual endpoint lengths and the full finite-white list}
\label{app:arithA17endpointlist}

Even with a golden end hole, the chord reference has energy
\[
 E_{\min}={13079276\over15000000}>.8719 .
\]
Each signed end block has length \(L=5\) or \(6\) and mean
\(.25<b<.3\).
If its actual endpoint coordinate were at least \(.63\), replace
the block by
\[
 .63,\quad (Lb-.63)/(L-1),\ldots,(Lb-.63)/(L-1).
\]
Its sum and every available central height are unchanged.
The remaining coordinates exceed \(.155\), and the following
coordinate is \(<.1\); hence the changed positive central
coefficient is \(>.055/\sqrt2>.018\).
All other positive central coefficients still meet \(.018\).
Thus the SAME summed-foot proof applies to this ONE new
reference, whose energy increases by
\[
 {L\over L-1}(.63-b)^2>{6\over5}(.33)^2=.13068.
\]
This exceeds one after adding \(E_{\min}\).
Therefore both signed actual endpoint coordinates are \(<.63\).

For either actual endpoint edge,
\[
 x_{\rm end}\ge\sqrt{17/18}
                        \sqrt{(C-1)/(C+1)} .
\]
The bound \(C\ge5/2\) would give endpoint square
\(\ge17/42>.63^2\), impossible.
If the ACTUAL field is nongolden, no central golden hole exists.
Use the full reference and graft instead to \(.575=23/40\).
The remaining block coordinates exceed \(.225125\);
all changed central coefficients still exceed \(.018\).
The energy gain is
\((5/4)(.575-.2952)^2=.09786005\), giving total
\(>1.00164385\).
Since \(C\ge21/10\) would give endpoint square
\(\ge187/558>(23/40)^2\), we obtain
\[
 \text{always }C<5/2,\qquad
 \text{outside }\mathbb Q(\varphi): C<21/10 .
\tag{A17C}\label{eq:arithA17C}
\]

Release old node one.
The retained white path is \(A_{16}\) on nodes two through
seventeen.  Finite one-node classification gives isolation;
helper three at \(j=2,3,16,17\); or helper four at \(j=2,17\).
A second contact would make a cycle; no exceptional finite
rank-seventeen diagram contains that path.
There is no foreign contact in this whole-diagram case.
Parallel own always has its actual future-null pair.

For helper three put \(t=2\gamma\).
Connected simple-three normalization gives a square-two cap,
helper product \(-1\), own product \(-t\), and integral \(t\).
Its full Lorentz star is defined over \(\mathbb Q(t)\);
all-place definiteness gives ACTUAL \(K=\mathbb Q(t)\).
The unit time square and edge coshes are exactly
\[
 A^2={17t^2+2(18-j)t+j(18-j)-36\over36},
\qquad
\begin{array}{c|c}
j&C\\ \hline
2&(17\gamma+8)/9\\
3&(17\gamma+15/2)/(3\sqrt2)\\
16&(17\gamma+1)/(3\sqrt2)\\
17&(17t+1)/18 .
\end{array}
\]
These are obtained by clearing ALL retained \(A_{16}\) roots.
For helper four the two formulas are
\[
 j=2:\ C=(17\gamma+8\sqrt2)/3,\qquad
 j=17:\ C=(17\gamma+1/\sqrt2)/3 .
\]
The raw white residuals are \(18/17\) at a helper-three
endpoint, \(4/17\) at a next-to-end three, and \(1/17\)
at a helper-four endpoint.

At helper \(j=3\), every positive own violates \(C<5/2\).
Own zero has the ACTUAL feasible affine-\(E_8\) null
\[
 2n_1+4n_2+6n_3+5n_4+4n_5+3n_6+2n_7+n_8+3c .
\]
Its future time is \(3A>0\), and all omitted products are
nonpositive.
At \(j=16\), own zero has \(A^2=-1/9\), own three has
the actual all-three cycle, and own at least four violates
\(C<5/2\).
The helper-four rows fail already at own zero for \(j=2\)
and own at least three for \(j=17\); the latter own-zero
time is \(-1/18\).
At \(j=2\), own zero has negative time and
\(C<5/2\) gives \(\gamma<29/34\), leaving owns three,
four or five.
Own three is an actual affine triangle.
Own four generates \(\mathbb Q(\sqrt2)\), but its edge exceeds
\(21/10\), contrary to the nongolden bound.
Only own five in the ACTUAL golden field can remain at this end.

\subsection{The far-helper strip and the golden exceptions}
\label{app:arithA17goldenends}

For the far helper \(j=17\), positive time gives \(t>1\).
Outside the golden field, \eqref{eq:arithA17C} gives
\[
 1<t<184/85,\qquad -19/17<t_\sigma<1\quad(\sigma\ne\mathrm{id}).
\]
Set \(f(s)=s(s^2-1)\).
On this nonphysical interval
\(\lvert f(s)\rvert\le2/(3\sqrt3)\):
on \([-1,1]\) use
\(4/27-q(1-q)^2=(1-3q)^2(4-3q)/27\);
on \([-19/17,-1]\) the bound \(1368/4913\) is smaller.
Physically \(0<f(t)<8\).
Degree at least four would give nonzero integer norm of
magnitude \(<8[2/(3\sqrt3)]^3<1\).
At degree three the magnitude is \(<32/27<2\), so
\(t,t-1,t+1\) would all be units.
The signs at one and minus one leave exactly
\[
 X^3-X-1,\quad X^3-2X^2-X+1,\quad
 X^3+X^2-2X-1,\quad X^3-X^2-2X+1 .
\]
The first has discriminant \(-23\), so is not totally real.
The second is negative at \(11/5\), with increasing positive
branch there, so its physical root exceeds \(11/5>184/85\).
The latter two have positive values \(5347/4913,2899/4913\)
at \(-19/17\), contrary to all three roots lying above that
endpoint.

In degree two the integral trace is \(0,1,2,3\), and the nonzero
norm is \(\pm1,\pm2\).
For trace zero the negative conjugate strip forces
\(1<t<19/17\), so \(-t^2\) lies strictly between \(-2,-1\).
For trace one the norm \(t(1-t)\) leaves only \(-1\), giving
\(t=\varphi\), or \(-2\), giving the rational \(t=2\).
For trace two the norm \(t(2-t)\) is strictly between \(-1,1\),
with zero excluded in degree two.
For trace three the conjugate strip requires \(t>2\), and
\(t(3-t)\) lies strictly between one and two.
Thus no nongolden quadratic survives.
In degree one only \(t=2\) is possible, the forbidden parallel.
Hence any far helper forces the actual golden field.

In \(K=\mathbb Q(\varphi)\), the unconditional bound gives
\(1<t<44/17\) and the same nonphysical strip.
Writing \(t=a+b\varphi\) in the FULL ring yields \(0<b<2\);
thus \(b=1\), \(a=0\) or one, and \(a=1\) violates the
physical bound.  Again \(t=\varphi\).
Every possible nonisolated endpoint is therefore a near- or
far-simple-three helper with own five in \(K=\mathbb Q(\varphi)\).

We now exclude that actual field in full, retaining its isolated
alternatives.  The unit near/far helper times are \(M_N/36,M_F/36\),
where
\[
 M_N=13+49\varphi,\quad \mathrm N(M_N)=-1595,\qquad
 M_F=-2+19\varphi,\quad \mathrm N(M_F)=-395 .
\]
Two far helpers have positive actual product \((1+4\varphi)/9\).
Near/far time-norm ratio \(319/79\) is not a rational square.
Two near helpers have invariant two-cycle
\((500+660\varphi)/81\notin\mathbb Z[\varphi]\).
Thus both endpoints cannot be nonisolated.

For an isolated endpoint put \(U=4\gamma^2\).
Rescale its actual \(K\)-root line to raw square \(2U\) and
sole old product \(-U\).
The actual two-cycle makes \(U\) integral and the full Lorentz
star gives \(K=\mathbb Q(U)\).
Its raw time square is
\[
 H^2=U(17U-36)/18,\quad
 36/17<U<225/17,\quad0<U'<36/17 .
\]
The FULL ring models \(U=a+b\varphi\) are exactly
\[
\begin{array}{c|r|r|r}
U&\mathrm N(U)&-\mathrm N[U(17U-36)]&
 \text{time norm squareclass}\\ \hline
1+\varphi&1&251&251\\
2+\varphi&5&1595&5\cdot11\cdot29\\
2+2\varphi&4&4880&5\cdot61\\
3+2\varphi&11&4631&11\cdot421\\
2+3\varphi&1&2699&2699\\
3+3\varphi&9&14499&179\\
3+4\varphi&5&16895&5\cdot31\cdot109\\
4+4\varphi&16&22784&89\\
4+5\varphi&11&38291&11\\
5+5\varphi&25&16475&659 .
\end{array}
\]
Indeed \(b\sqrt5<225/17\) gives \(b=1,\ldots,5\);
the strict conjugate strip gives respectively
\(a=1,2;\ 2,3;\ 2,3;\ 3,4;\ 4,5\).
The entries follow from
\[
 \mathrm N[U(17U-36)]
 =n(289n-612T+1296),\quad n=\mathrm N(U),\quad T=\mathrm{Tr}(U).
\]
The displayed squareclasses are distinct; their small prime
factorizations can be checked by trial divisibility through the
respective square roots.

Actual roots on the SAME \(K\)-time line require
\(U(17U-36)/M_N\) or \(/M_F\) to be a square IN \(K\).
Their rational norms exclude every model except \(U=2+\varphi\)
with a near helper; none has the far-helper class \(5\cdot79\).
For that exception \(U(17U-36)=M_N\) exactly.
Its equal positive raw time coefficients and spatial product
\(U^2/18\) give
\[
 g=-(4+22\varphi)/9,\qquad
 g^2/U=(168+164\varphi)/81,
\]
whose nonintegral two-cycle excludes it.

Finally two isolated golden endpoints cannot use different rows,
since their rational time-norm squareclasses differ.
With the same row their equal positive raw times give
\[
 g=2U(9-4U)/9,\qquad g^2/U^2=4(9-4U)^2/81\in\mathcal O_K .
\]
Integral closedness makes \(2(9-4U)/9\) integral, hence
\(8U/9\) integral; Bezout \(8\cdot8-7\cdot9=1\)
then makes \(U/9\) integral.
But \(0<\mathrm N(U)<8100/289<29\), whereas
\(\mathrm N(U/9)=\mathrm N(U)/81\) would be a positive integer.
Thus the actual golden field is impossible.

Both endpoints are now isolated in a nongolden actual field,
with independent integral generators satisfying
\[
 K=\mathbb Q(U)=\mathbb Q(V),\quad
 36/17<U,V<3969/425,\quad
 0<U^\sigma,V^\sigma<36/17\quad(\sigma\ne\mathrm{id}).
\tag{A17isolated}\label{eq:arithA17isolated}
\]
No equality of their parameters, finite-own assumption or common
disconnected raw scale was used.

\subsection{One two-end graft gives actual degree at most five}
\label{app:arithA17degree}

There is no golden hole after the preceding argument.
If BOTH signed endpoint coordinates were at least \(1/2\),
modify the first and last five-coordinate blocks simultaneously.
Replace each signed block mean \(b\) by
\[
             a,r,r,r,r,\quad a=1/2,\quad r=(5b-a)/4 .
\]
Both sums and every central height stay unchanged.
The supports are disjoint and each residual has sum zero,
so the two residuals are orthogonal to one another and to the
old flat blocks.
Moreover \(r>.243875>.1\), making every changed positive
central coefficient \(>.018\).
The new endpoint coefficients are paid by the two ACTUAL
coordinates assumed at least \(a\).
The summed-foot proof therefore applies to this ONE new reference,
with energy
\[
 >.9037838+2(5/4)(.5-.2952)^2
 =1.0086414 .
\]
Thus at least one actual endpoint coordinate is \(<1/2\).
For its generator \(U\), the edge formula
\(C^2=17U/36\) and the farthest endpoint comparison give
\[
                  36/17<U<66564/10625<157/25.
\tag{A17shortU}\label{eq:arithA17shortU}
\]
The other generator need not satisfy this sharper upper bound;
the chosen one still generates the ENTIRE actual field.

Let \(F(U)=U(U-1)(U-2)\) and \(d=[K:\mathbb Q]\).
For \(d\ge2\) all factors are nonzero algebraic integers.
Physically
\[
                    0<F(U)<142 .
\]
At any other embedding put \(t=U^\sigma-1\in(-1,19/17)\).
On \([-1,1]\), \(|t^3-t|\le2/(3\sqrt3)<2/5\);
on \((1,19/17)\) it is at most \(1368/4913<1/3\).
If \(d\ge7\), the nonzero integral norm would have magnitude
\(<142(2/5)^6=9088/15625<1\).
Hence \(d\le6\).
For \(d=6\), the bound \(142(2/5)^5<2\) makes its absolute
norm one, and all three separate factor norms have absolute
value one.

We give the complete degree-six seal.
Set \(s=U-1\), and let
\(r_i=|U^{\sigma_i}-1|\), \(i=1,\ldots,5\).
Irreducibility excludes \(r_i=0,1\).
The unit norms and~\eqref{eq:arithA17shortU} give
\[
 0<r_i<19/17<9/8,\quad
 \prod r_i=1/s>3/16,\quad s<132/25,\quad
 s^2-1<16799/625<27.
\]
The sharper \(s<132/25\), rather than merely \(s<16/3\),
is needed in the last inequality.

If \(m\) positive numbers \(q_i<1\) have
\(p=(\prod q_i)^{2/m}\), ordinary AM--GM applied separately
to \(q_i^2\) and \(1-q_i^2\) gives
\[
       p+\{\prod(1-q_i^2)\}^{1/m}\le1 .
\]
Let \(k\) be the number of \(r_i>1\), so \(m=5-k\).
Each large \(r_i\) contributes
\(|1-r_i^2|<72/289<1/4\).
The following six cases exhaust all possibilities:
\[
\begin{array}{c|c|c|c}
k&\prod_{\rm small}r_i\text{ lower}&
 p\text{ lower}&\prod_{i=1}^5|1-r_i^2|\text{ upper}\\ \hline
0&3/16&1/2&1/32\\
1&1/6&2/5&81/2500\\
2&4/27&1/4&27/1024\\
3&32/243&1/8&49/4096\\
4&\text{not needed}&\text{not needed}&1/256\\
5&\text{empty product}&\text{not needed}&1/1024 .
\end{array}
\]
For the product bounds use
\(\prod_{\rm small}r_i>(3/16)(8/9)^k\).
The four root comparisons for \(p\) follow respectively from
\(9/256>1/32\), \(1/6>4/25\), \(16/729>1/64\),
\(32/243>1/8\); each exponent has been cleared positively.
Multiplying the last column by \(27\) gives
\[
 27/32,\quad2187/2500,\quad729/1024,\quad
 1323/4096,\quad27/256,\quad27/1024,
\]
all below one.
But the unit norm of \(U(U-2)\) would be
\[
             1=(s^2-1)\prod_{i=1}^5|1-r_i^2|<1 .
\]
This excludes degree six and proves the ACTUAL degree bound
\(d\le5\), without a field list or an equality assumption on the
two endpoints.

\subsection{The whole two-dotted interval and the derived finite fields}
\label{app:arithA17anchors}

Suppose both isolated owns are dotted.
For the left endpoint put \(u=\gamma^{-2}\in(0,1)\).
The positive actual core
\[
 c+\gamma^{-1}n_1
   =\gamma\{\sqrt{17/9-u}\,p-(\omega_1-un_1)\}
\]
has negative square \(u-1\), zero \(c\) product and negative
\(n_1\) product.
Every other original wall, INCLUDING the other endpoint opposite,
has nonpositive product with both summands.
After normalization it is a genuine chamber point and gives
\[
 x\cdot Z_L\ge h(u)=\sqrt{17/9-u}-\sqrt{u(1-u)},\quad
 Z_L=\omega_1-un_1 .
\]
The right endpoint has its independent parameter \(v\in(0,1)\)
and reflected vector \(Z_R\).

Let
\[
 Y_L=Z_L-\omega_5/5,\qquad
 R(u)=\|Y_L\|^2=8/5-2u+u^2 .
\]
This is a zero-sum vector supported in the first five coordinates.
Its central-root products are zero except \(-1/5\) at five.
Its reflected partner is supported in the LAST five coordinates.
Consequently \(w\cdot Y_L=w\cdot Y_R=Y_L\cdot Y_R=0\), while
\(w\cdot Z_L=w\cdot Z_R=.4174\).
Pay both actual core pressures at \(\mu=.3\):
\[
 w'=w+\mu Y_L+\mu Y_R,\quad
 \beta'_5=\beta'_{13}=.1382-.06=.0782,\quad
 \beta'_i=\beta_i\ (6\le i\le12).
\]
Every positive central coefficient is still at least \(.018\),
and \(w'\cdot n_i=\beta'_i\) EXACTLY.

Retain every central isolated foot, with its OWN parameter
\(v_i\in(0,1)\), using ONE vector
\[
                       V=w'-\sum_{i\in S}\beta'_iv_in_i .
\]
Sum the endpoint and central actual pressures before estimating
the norm.  The identities just proved and
\(n_i\cdot n_j\le0\) give
\[
 \begin{split}
 x\cdot V&\ge E_0+
 \mu\{h(u)-.4174+h(v)-.4174\}
       -\sum_S\beta'_i(\lambda_i-.6),\\
 \|V\|^2&\le E_0+\mu^2\{R(u)+R(v)\}
       -\sum_S(\beta'_i)^2(2v_i-v_i^2).
 \end{split}
\]
Every central summand is nonnegative by the already proved
absorption inequality, since \(\beta'_i\ge.018\).
Thus
\[
 \|x\|^2\ge E_0+G(u)+G(v),\qquad
 G(u)=.6\{h(u)-.4174\}-.09(8/5-2u+u^2).
\]

This whole continuous interval has a strict rational seal.
Concavity and the endpoint squares give
\(\sqrt{17/9-u}>1.3743-.4315u\) on \([0,1]\).
Therefore
\[
 G(u)>.43014-.0789u-.09u^2-.6\sqrt{u(1-u)}.
\]
For \(u\le1/2\), use \(u^2\le u/2\) and the circle inequality
\[
 au+b\sqrt{u(1-u)}\le(a+\sqrt{a^2+b^2})/2 .
\]
With \(a=.1239,b=.6\), \(\sqrt{a^2+b^2}<.6127\)
gives \(G(u)>.06184\).
For \(u\ge1/2\), use \(u^2\le1.5u-.5\);
\(a=.2139,b=.6\), with radical \(<.637\), gives
\(G(u)>.04969\).
The positive square margins are \(.00005008,.00001579\),
respectively.
Hence throughout \((0,1)\),
\[
          \|x\|^2>.9037838+2(.04969)=1.0031638,
\]
impossible.  Parallel owns already have their actual future null.
At least one endpoint is a finite native own.

At that endpoint \(U=4\cos^2(\pi/m)=2+2\cos(2\pi/m)\)
generates the ACTUAL field.
The degree bound makes \(\phi(m)\le10\).
The COMPLETE elementary totient list is
\[
 3,4,5,6,7,8,9,10,11,12,14,15,16,18,20,22,24,30 .
\]
Indeed only prime divisors \(2,3,5,7,11\) are possible;
their permissible prime powers and products give that list.
Labels three and four have \(U\le2<36/17\).
Six gives the rational field, excluded by Meyer and compact
anisotropy at form rank eighteen.
Five and ten give the already excluded golden field.
For the following nine labels choose the indicated coprime
cyclotomic embedding:
\[
\begin{array}{c|c|c|c}
m&k&\cos(2k\pi/m)=\sin\theta&2\theta/\pi\\ \hline
9&2&\theta=\pi/18&1/9\\
14&3&\theta=\pi/14&1/7\\
15&2&\theta=7\pi/30&7/15\\
16&3&\theta=\pi/8&1/4\\
20&3&\theta=\pi/5&2/5\\
24&5&\theta=\pi/12&1/6\\
30&7&\theta=\pi/30&1/15\\
11&2&\theta=3\pi/22&3/11\\
22&3&\theta=5\pi/22&5/11 .
\end{array}
\]
Every last entry exceeds \(1/17\).
Strict concavity on \(0<\theta<\pi/2\) gives
\(\sin\theta>2\theta/\pi\), so that conjugate would have
\(U^\sigma>36/17\), a contradiction.
Label thirty was included before its conjugate test.
Only labels \(7,8,12,18\) remain.
Their actual fields are, respectively,
\(K_7,\mathbb Q(\sqrt2),\mathbb Q(\sqrt3),K_9\);
the cubic parameters have polynomials
\[
 X^3-5X^2+6X-1,\qquad X^3-6X^2+9X-3,
\]
with discriminants \(49,81\) and irreducibility by the
rational-root test.  This proves the actual field reduction,
rather than assuming a list of possible fields.

\subsection{Actual inclusion in the recorded finite domains}
\label{app:arithA17datainclusion}

The isolated endpoint normalization in
\eqref{eq:arithA17isolated} has raw square \(2U\), sole old
product \(-U\) and time square \(F(U)/18\), where now
\(F(X)=X(17X-36)\).
For the two actual endpoint roots, subtracting their \(K\)-valued
old projections leaves one common \(K\)-time line.
Thus
\[
 R=F(U)F(V)\in K^2,\quad
 g={UV-\sqrt R\over18},\quad
 L={g^2\over UV}.
\]
The square root is PHYSICALLY POSITIVE.
The long old-three cycle gives \(g\in\mathcal O_K\);
the actual two-cycle gives \(L\in\mathcal O_K\).
Acuteness gives \(g\le0\).
If \(g\ne0\), every nonphysical place satisfies \(0<L^\sigma<4\);
the branch \(g=L=0\) is retained.
The physical owns and mutual products retain their native/dotted
alphabet, with a parallel pair excluded by its actual future null.
These are precisely the stated necessary tests in the quadratic
domain; the endpoint parameters need not be equal.

If \(K\) is quadratic, this actual pair therefore enters the FULL
quadratic domain of Proposition~\ref{prop:pair-output}(i),
including its maximal-order classes, square signs and zero mutual.
That recorded finite-output proposition says the domain is empty.
If \(K=K_7\) or \(K_9\), the finite endpoint proved above supplies
the actual anchor seven or eighteen, respectively.
The same original time-line square, cyclic and native conditions
put the other endpoint in the full cubic finite-anchor domain
of Proposition~\ref{prop:pair-output}(ii).
Its complete ring/box/decision arguments are given in
Section~\ref{sec:algorithms}, and the recorded output is empty.

Thus, at those explicitly stated EXISTING finite-output premises,
the whole \(A_{17}\) at the original arithmetic rank-seventeen
vertex is impossible.
The geometric domain inclusion and complete hand arguments are
proved in this appendix; no new enumeration, independent replay,
runtime PASS, higher-dimensional arithmetic tail or geometric
Lean formalization is inferred by transcribing them.

\section{The separate whole-\(A_{16}\) endpoint proof}
\label{app:wholearithA16}

The whole \(A_{16}\) case has \(N=17\) and its own constants.
It is not a consequence of reranking the preceding \(A_{17}\)
theorem.  We retain one original farthest vector throughout.
The common chain proof gives central reference energy
\[
 E_0=.8708972,\quad
 (H_5,H_6,H_7,H_8)=(2.056,2.186,2.269,2.310),\quad H_9=H_8,
\]
with positive coefficients at least \(.0205\).
Its worst single end-hole chord has energy
\(12569249/15000000\).
All grafts below preserve the common simultaneous-foot proof
by checking those coefficients afresh.

\subsection{Actual lengths and all nonisolated endpoint roles}
\label{app:arithA16endpointlist}

On either signed end block, of length five or six and mean
\(.25<b<.3\), grafting to \(.67\) leaves each remaining
coordinate \(>.145\); the following coordinate is below \(.1\).
The central coefficient stays \(>.045/\sqrt2>.0205\).
The gain \(>(6/5)(.37)^2=.16428\), even after the worst
chord omission, exceeds the available norm budget.
Both signed endpoint coordinates are therefore \(<.67\).
The actual edge comparison is now
\[
 x_{\rm end}\ge\sqrt{16/17}\sqrt{(C-1)/(C+1)}.
\]
At \(C=17/6\) its square \(176/391\) exceeds \(.67^2\).
Hence both endpoints have \(C<17/6\).
If both signed coordinates were at least \(.56\), disjoint
grafts on the SAME reference would gain
\(>2(6/5)(.26)^2=.16224\), again exceeding one.
Since \(16/51>.56^2\), at least one endpoint has \(C<2\).

Outside the actual golden field there is no hole and the
full block mean lies in \((.2907,.2908)\).
A single graft to \(.613\) gains \(>.12976605\), making total
\(>1.00066325\).
The retained coordinates are \(>.210125\), so the coefficient
threshold remains valid.
As \(32/85>.613^2\), both endpoints then have \(C<7/3\).
Two disjoint grafts to \(.52\) gain \(>.1313316\), giving total
\(>1.0022288\).
The exact comparison
\(1296/4777>(13/25)^2\) therefore proves
\[
 \text{always both }C<17/6,\ \min C<2;\qquad
 K\ne\mathbb Q(\varphi):
        \text{ both }C<7/3,\ \min C<181/100 .
\tag{A16lengths}\label{eq:arithA16lengths}
\]

Release node one and retain the whole white \(A_{15}\).
Its complete finite one-node extensions are isolation;
helper three at \(j=2,3,15,16\); or helper four at \(j=2,16\).
There is no foreign component or exceptional rank-sixteen
completion.  A second contact makes a cycle.
For helper three, the actual connected normalization gives
integral \(t=2\gamma\), ACTUAL \(K=\mathbb Q(t)\), raw
time square \(M_j/17\), and
\[
 M_j=16t^2+2(17-j)t+j(17-j)-34,\qquad
\begin{array}{c|c}
j&C\\ \hline
2&(16t+15)/17\\
3&(16\gamma+7)/\sqrt{17}\\
15&(16\gamma+1)/\sqrt{17}\\
16&(16t+1)/17 .
\end{array}
\]
For helper four the lengths are
\[
 (16\gamma+15/\sqrt2)/\sqrt{17/2},\qquad
 (16\gamma+1/\sqrt2)/\sqrt{17/2}.
\]
These follow by eliminating ALL retained roots; the unit white
residuals are \(17/32,1/8,1/16\) in the endpoint-three,
next-three and endpoint-four rows.

Helper \(j=3\) with positive own gives
\(C\ge15/\sqrt{17}>17/6\).
Its own-zero case has the same actual affine-\(E_8\) vector
printed in the preceding endpoint proof.
At \(j=15\), own zero has unit time square \(-2/17\), own three
has the actual all-three cycle, and owns at least four give
\(C>17/6\).
Helper four at \(j=2\) fails even at own zero.
At \(j=16\), its own-zero time is \(-1/17\), and every positive
own gives \(C\ge(8\sqrt2+1)/\sqrt{17}>17/6\).
At near \(j=2\), own zero has negative time and own three
the actual affine triangle.
At far \(j=16\), positive time requires \(t>1\), since
\(M_{16}=(t-1)(16t+18)\).
Parallel has the positive actual future-null pair.
All these nulls use UNIT actual normals and are feasible
against every omitted wall.

In a nongolden field the near helper has \(t<37/24\),
so its only remaining native own is four, \(t=\sqrt2\).
The far helper satisfies
\[
                    1<t<29/12,\qquad -9/8<t_\sigma<1.
\]
As before \(f(t)=t(t^2-1)\) has other-place magnitude at most
\(2/(3\sqrt3)\); on the additional interval
\([-9/8,-1]\), the maximum is \(153/512\), smaller.
Physically \(0<f(t)<12\).
The integral norm excludes degrees at least four.
At degree three \(t,t-1,t+1\) would all be units.
The same complete four cubics as before occur.
The negative-discriminant cubic and the two with a conjugate
below \(-9/8\) fail; the only survivor is
\[
                P(X)=X^3-2X^2-X+1,
\]
whose physical root \(a>2\) generates \(K_7\).
For degree two, trace \(0,1,2,3\) and norm
\(-2,-1,1,2\) exhaust the strip.
Trace zero would require \(\sqrt2\), with conjugate below
\(-9/8\); trace one leaves only \(\varphi\) or rational two;
trace two leaves \(1+\sqrt2\), since \(1+\sqrt3\) is too large;
trace three's irrational positive root exceeds \(29/12\).
The rational two is parallel.
Thus the nongolden helper roles are exactly near silver,
far silver, or far \(K_7\).
Their actual lengths are respectively
\[
 (16\sqrt2+15)/17,\quad (17+16\sqrt2)/17,\quad
 (16a+1)/17,
\]
all \(>181/100\).
The short endpoint selected in~\eqref{eq:arithA16lengths}
is consequently ISOLATED.

\subsection{The full golden endpoint table}
\label{app:arithA16golden}

For any isolated endpoint, its OWN \(K\)-rescaling gives
raw norm \(2U\), sole old product \(-U\), integral \(U=4\gamma^2\)
generating ACTUAL \(K\), and
\[
 F(U)=U(16U-34),\quad H_U^2=F(U)/17,\quad C^2=8U/17,
 \quad U>17/8,\quad0<U^\sigma<17/8 .
\tag{A16U}\label{eq:arithA16U}
\]
No raw time itself is assumed \(K\)-valued; the ratios along
the SAME one-dimensional \(K\)-time line are \(K\)-valued.

Suppose \(K=\mathbb Q(\varphi)\).
The near helper has \(t<199/96\) and conjugate
\(-2<t'<1/8\).  The finite native list leaves \(\varphi\);
a dotted full-ring \(a+b\varphi\) has \(b=1\) and no integer
\(a\) in \(2<t<199/96\).
The far helper satisfies \(1<t<283/96\), \(-9/8<t'<1\),
giving \(t=\varphi\) or \(\varphi^2\).
All three types have raw times \(M/17\):
\[
\begin{array}{c|c|r}
\text{helper}&M&\mathrm N(M)\\ \hline
\text{near }\varphi&12+46\varphi&-1420\\
\text{far }\varphi&-2+18\varphi&-356\\
\text{far }\varphi^2&16+50\varphi&-1444 .
\end{array}
\]
The norm squareclasses \(355,89,1\) forbid different types.
Reflected near copies have nonintegral actual long-cycle
\(g=-(7+41\varphi)/17\).
Reflected far-\(\varphi\) copies have positive product
\(g=(4+15\varphi)/17\).
Reflected far-\(\varphi^2\) copies have
\(g=(19-15\varphi)/17\in(-1,0)\), forbidden for two
distinct square-two actual roots by the native gap.
Thus both endpoints cannot be nonisolated.

For the isolated golden end, the universal length cap gives
\[
                       17/8<U<4913/288,\qquad0<U'<17/8 .
\]
Writing \(U=a+b\varphi\) gives \(b=1,\ldots,7\).
The complete conjugate intervals give
\(a=1,2;\ 2,3;\ 2,3;\ 3,4;\ 4,5;\ 4,5;\ 5,6\),
and the physical upper bound removes only \((6,7)\).
The complete thirteen rows are
\[
\begin{array}{c|r|r|r}
U&\mathrm N(U)&-\mathrm N(F(U))&\text{squareclass}\\ \hline
1+\varphi&1&220&55\\
2+\varphi&5&1420&355\\
2+2\varphi&4&4336&271\\
3+2\varphi&11&4180&1045\\
2+3\varphi&1&2396&599\\
3+3\varphi&9&12924&359\\
3+4\varphi&5&15020&3755\\
4+4\varphi&16&20416&319\\
4+5\varphi&11&34100&341\\
5+5\varphi&25&15100&151\\
4+6\varphi&4&21744&151\\
5+6\varphi&19&50996&12749\\
5+7\varphi&11&58036&14509 .
\end{array}
\]
For \(n=\mathrm N(U),T=\mathrm{Tr}(U)\), every entry follows from
\(\mathrm N(F(U))=n(256n-544T+1156)\).
The listed classes are squarefree:
\(1045=5\cdot11\cdot19\), \(3755=5\cdot751\),
\(319=11\cdot29\), \(341=11\cdot31\),
\(12749=11\cdot19\cdot61\), \(14509=11\cdot1319\),
\(355=5\cdot71\), \(55=5\cdot11\);
the remaining small primality checks are direct trial division.
The only repeated class is151.
Precisely the FIRST SIX rows have \(U<17/2\), equivalently \(C<2\),
and their classes occur nowhere else in this complete table.

If one endpoint is near-\(\varphi\) or far-\(\varphi^2\),
its \(C>2\), so the opposite isolated end is in the first six rows.
Class one matches none.
The near class355 matches only \(U=2+\varphi\);
there \(F(U)=12+46\varphi\), so the actual positive raw
times are equal.
Their spatial product is \(U(\varphi+2)/17=U^2/17\),
giving the same nonintegral long cycle
\(g=-(7+41\varphi)/17\).
For a far-\(\varphi\) helper the opposite end may be ANY of
the thirteen rows, but none has class89.
All helper/isolated possibilities are excluded.

Both golden endpoints must now be isolated.
One belongs to the first six rows by the selected \(C<2\).
The SAME time-line square test forces equality of the two
\(U\)-models; the repeated151 class cannot occur at the selected end.
Equal positive raw times give
\[
 g=U(34-15U)/17,\qquad
 L=g^2/U^2=((34-15U)/17)^2\in\mathcal O_K .
\]
Integral closedness makes \((34-15U)/17\), and then \(U/17\),
integral, using Bezout for15 and17.
But the selected norms \(1,5,4,11,1,9\) are smaller than289,
so \(\mathrm N(U/17)=\mathrm N(U)/289\) cannot be a
nonzero integer.
The ACTUAL golden field is impossible.

\subsection{Both actual endpoints are isolated}
\label{app:arithA16helpers}

The actual field is now nongolden, and the selected endpoint
of \(C<181/100\) is isolated.
At \(K=\mathbb Q(\sqrt2)\), its full possible isolated strip
from \(C<7/3\) is
\[
 17/8<U<833/72,\quad0<U'<17/8 .
\]
Writing \(U=A+B\sqrt2\), the complete models are
\[
 2+s,\ 3+s,\ 3+2s,\ 4+2s,\ 5+3s,\ 6+3s,\qquad s=\sqrt2 .
\]
Indeed \(B=1,\ldots,4\);
the conjugate strip permits respectively
\(A=2,3;\ 3,4;\ 5,6;\ 6,7\), and the physical upper
bound removes both \(B=4\) possibilities.
Their signed \(F\)-norms are
\[
                   -1016,-2212,-1852,-9184,-17444,-13752 .
\]
The near-silver helper has \(M=28+30s\), norm
\(-1016=-8\cdot127\).
Only \(U=2+s\) has a compatible norm squareclass:
none of the other five norms is divisible by127.
For that exception \(F(U)=M\) and the actual long cycle
\(g=-(22+26s)/17\) is not integral.
The far-silver helper has \(M=32+34s\), norm
\(-1288=-8\cdot7\cdot23\).
None of the six time norms is divisible by23.
Both silver helper roles are excluded.

For the remaining far helper in \(K_7\), let \(a>2\) be
the root of \(P(X)=X^3-2X^2-X+1\).
Its three roots have disjoint rational brackets
\[
 2.246<a<2.247,\quad -.803<b<-.801,\quad .554<c<.556 .
\]
Signs and derivatives of \(P\) isolate them.
The full ring is \(\mathbb Z[a]\), as justified below.
The selected isolated endpoint has
\(17/8<U<7\), \(0<U_b,U_c<17/8\).
Write \(U=A+Ba+Ca^2\).
Lagrange interpolation expresses \(C\) as the sum of the
three values divided by the derivatives at the three roots.
The two positive denominators exceed five and four;
the magnitude of the negative denominator exceeds \(17/8\).
Consequently \(-1<C<7/5+(17/8)/4<2\), so \(C=0,1\).
The small conjugate difference gives
\[
                 |B+C(2-a)|<(17/8)/(c-b)<8/5 .
\]
If \(C=0\), positive physical difference forces \(B=1\),
and the conjugate interval forces \(A=1\).
If \(C=1\), the physical/conjugate difference \(<7\)
removes \(B\ge1\), leaving \(-1,0\).
For \(B=-1\), the \(c\)-value requires \(A\ge1\), while
the \(b\)-value requires \(A\le0\).
For \(B=0\), \(A=0,1\).
Thus the COMPLETE selected models are
\[
                               U=1+a,\quad a^2,\quad1+a^2 .
\]
Their positive time norms are \(5048,32248,82264\).
The helper time numerator is
\((a-1)(16a+18)\), of norm \(7496=8\cdot937\).
The prime937 divides none of these three norms;
therefore no actual time ratio can be a \(K\)-square.
This excludes the last nongolden helper.
BOTH actual endpoints are isolated.

\subsection{Degree, the whole two-dotted domain and finite anchors}
\label{app:arithA16anchors}

The selected isolated endpoint generates the actual field, with
\[
 17/8<U<556937/80000<7,\qquad0<U^\sigma<17/8 .
\]
For \(F_3(U)=U(U-1)(U-2)\), physical value \(0<F_3(U)<210\).
At a nonphysical embedding \(y=U^\sigma-1\in(-1,9/8)\).
The same cubic bound gives
\(|y^3-y|<2/5\), including the upper extra interval where
the maximum is \(153/512\).
A degree at least seven would give a nonzero integer norm
\(<210(2/5)^6=.86016<1\).
Hence ACTUAL degree is at most SIX.
No degree-five refinement from the \(A_{17}\) case is used.

Suppose both isolated owns were dotted.
Their ACTUAL endpoint cores give, for independent \(u,v\in(0,1)\),
\[
 x\cdot Z_L\ge h(u)=\sqrt{32/17-u}-\sqrt{u(1-u)},\qquad
 Z_L=\omega_1-un_1 ,
\]
and its reflection.
The core is future timelike, zero on its cap, negative on the
released old root, and nonpositive against EVERY other actual
wall, including the opposite endpoint.
The residual \(Y_L=Z_L-\omega_5/5\) has zero sum,
support in the first five coordinates, norm
\(R(u)=8/5-2u+u^2\), and sole central product \(-1/5\)
at five.
The reflected residual subtracts \(\omega_{12}/5\), and is
orthogonal to both the first residual and the old flat block.
Here \(w\cdot Z_L=w\cdot Z_R=.4112\).

Pay each core at weight \(\mu=.32\).
The individual outer coefficients change from \(.1406\) to
\(.1406-.064=.0766>.0205\);
every other central coefficient is unchanged.
In particular the changed reference has its EXACT products
equal to these new coefficients.
Subtract all retained central isolated-foot terms with their
own parameters, and apply the common simultaneous absorption
proof to this ONE vector.
All endpoint/central and adjacent-central crosses are retained,
giving
\[
 \|x\|^2\ge E_0+G(u)+G(v),\quad
 G(u)=.64\{h(u)-.4112\}-.1024(1.6-2u+u^2).
\]
The following COMPLETE three-chord certificate covers \([0,1]\):
\[
\begin{array}{c|c|c|c|c}
u\text{ interval}&L(u)<\sqrt{32/17-u}&K&a&
 \sqrt{a^2+.64^2}\text{ upper}\\ \hline
[0,1/2]&1.3719-.3924u&.451008&.097536&.6475\\
[1/2,3/4]&1.3989-.4464u&.506688&.208896&.6733\\
[3/4,1]&1.4385-.4992u&.570432&.293888&.7043 .
\end{array}
\]
The strict endpoint values \(1.3719,1.1757,1.0641,.9393\)
have squares below \(32/17,47/34,77/68,15/17\).
Concavity proves each chord lower bound.
On \([l,r]\) use \(u^2\le(l+r)u-lr\);
then \(G(u)>K-au-.64\sqrt{u(1-u)}\).
The circle inequality used earlier gives the three lower values
\(.07849,.06559,.071338\).
The squared radical margins are
\(.000142978704,.000095351184,.000068333456\).
Thus \(G(u)>.06559\) on the whole dotted interval, and
\[
                  \|x\|^2>.8708972+2(.06559)=1.0020772 .
\]
Both owns cannot be dotted.
Parallel has its actual future null, so at least one endpoint
is a finite native own.

Its integral \(U=2+2\cos(2\pi/m)\) generates the actual field,
and the degree bound gives \(\phi(m)\le12\).
The full totient list is
\[
 3,4,5,6,7,8,9,10,11,12,13,14,15,16,
 18,20,21,22,24,26,28,30,36,42 .
\]
Only prime factors \(2,3,5,7,11,13\) and their allowed prime
powers occur; checking products with totient at most twelve
gives exactly this list.
Labels three and four have \(U\le2<17/8\).
For own six the ACTUAL unit-root null is
\(\sqrt3c+2n_1+n_2\), a nonzero future ray feasible against
all other walls.
Five and ten give the excluded golden field.
The remaining forbidden labels have explicit coprime embeddings:
\[
\begin{array}{c|c|c}
m&k&\cos(2k\pi/m)\text{ lower}\\ \hline
9,11,13,15,21&2&\sin(\pi/18)>1/9\\
14,16,20,22,26,28&3&\sin(\pi/14)>1/7\\
24,36,42&5&\sin(\pi/12)>1/6\\
30&7&\sin(\pi/30)>1/15 .
\end{array}
\]
All exceed \(1/16\), by strict sine concavity.
Their conjugate \(U\) would exceed \(17/8\).
Only finite anchors seven, eight, twelve and eighteen remain,
with ACTUAL fields \(K_7,\mathbb Q(\sqrt2),\mathbb Q(\sqrt3),K_9\).

We use FULL integer rings.
For \(K_7\), \(P(X+3)=X^3+7X^2+14X+7\) is Eisenstein
at seven, and the polynomial discriminant is49.
For \(K_9\), put \(Q(X)=X^3-3X-1\);
\(Q(X+1)=X^3+3X^2-3\) is Eisenstein at three and the
discriminant is81.
The Eisenstein uniformizer criterion proves maximality at the
only possible index primes, while the discriminants exclude
every other index prime.
Equivalently these are the full-ring proofs already given in
Section~\ref{sec:algorithms}, under translation.
Thus the bases are the full \(\mathbb Z[a]\), \(\mathbb Z[k]\);
the two quadratic rings are \(\mathbb Z[\sqrt2]\) and
\(\mathbb Z[\sqrt3]\).

\subsection{All twenty-one actual finite-anchor time norms}
\label{app:arithA16times}

The opposite isolated endpoint \(V\), whether finite or dotted,
has
\[
                    17/8<V<833/72,\qquad0<V^\sigma<17/8 .
\]
For \(\mathbb Q(\sqrt2)\), its complete six models were printed
above.
For \(\mathbb Q(\sqrt3)\), the full ring gives
\(B=1,2,3\), and conjugate intervals
\(A=2,3;\ 4,5;\ 6,7\).
The physical upper bound removes only \((A,B)=(7,3)\).
The EXACT five models are
\[
                    2+r,\ 3+r,\ 4+2r,\ 5+2r,\ 6+3r,
                    \qquad r=\sqrt3.
\]

For \(K_7\), the leading-coefficient interpolation bound now
gives \(-1<C<3\), so \(C=0,1,2\) in \(V=A+Ba+Ca^2\).
The same small conjugate difference gives
\(|B+C(2-a)|<8/5\).
The complete integer possibilities are
\[
\begin{array}{c|c|c}
C&B&A\\ \hline
0&1&1\\
1&-1&\text{none}\\
1&0&0,1\\
1&1&1\\
2&-1,0&0 .
\end{array}
\]
For \(C=2,B\ge1\), the physical/conjugate difference exceeds
\(3.047(1+2(1.444))=11.846736>833/72\).
This excludes every remaining larger \(B\);
the conjugate-difference bound excludes the smaller ones.
The resulting EXACT six models are
\[
                 1+a,\ a^2,\ 1+a^2,\ 1+a+a^2,\ 2a^2-a,\ 2a^2 .
\]

For \(K_9\), the three roots have brackets
\[
 1.879<k<1.880,\quad -.348<b<-.347,\quad -1.533<c<-1.532 .
\]
The two positive Lagrange denominators exceed \(7.5,4\),
and the magnitude of the negative one exceeds \(17/8\).
Thus again \(C=0,1,2\).
The full conjugate difference gives \(|B-Ck|<9/5\).
The constant-term interval cases are
\[
\begin{array}{c|c|c}
C&B&A\\ \hline
0&1&2\\
1&1&1\\
1&2&1,2\\
1&3&3\text{ fails the physical upper bound}\\
2&2&\text{none}.
\end{array}
\]
For \(B=3,C=1\), the physical value is
\(>3+3(1.879)+1.879^2=12.167641>833/72\).
For \(C=2,B\ge3\), the physical difference exceeds
\(3.411(3+.694)>833/72\).
At \(C=2,B=2\), the \(b\)-strip forces \(A\ge1\) and
the \(c\)-strip \(A\le0\).
Hence the EXACT four models are
\[
                     2+k,\ 1+k+k^2,\ 1+2k+k^2,\ 2+2k+k^2 .
\]
Every listed accepted model satisfies all strict strips by the
root brackets; every omitted integer coefficient has been excluded.
These are complete handwritten domains, not unexecuted worker results.

The original raw times of two isolated roots are \(F(U)/17\),
\(F(V)/17\).
The same actual \(K\)-time line makes their ratio a \(K\)-square.
A necessary condition is a square rational norm; BOTH possible
signed norms of any putative square root satisfy that condition.
Here are ALL twenty-one signed time norms, in the stated model order:
\[
\begin{array}{c|l}
\text{field and finite anchor}&\mathrm N(F(V))\\ \hline
\sqrt2,\ 2+\sqrt2&
 -1016,-2212,-1852,-9184,-17444,-13752\\
\sqrt3,\ 2+\sqrt3&
 -764,-3432,-8688,-12428,-27612\\
K_7,\ 1+a&
 5048,32248,82264,110552,1784,330688\\
K_9,\ 2+k&
 16872,66168,79032,293624 .
\end{array}
\]
For a quadratic model the reproducible identity is
\[
 \mathrm N(F(V))=n(256n-544T+1156).
\]
For a cubic model with coefficient data \((T,S,n)\), it is
\[
 \mathrm N(F(V))
      =n(4096n-8704S+18496T-39304).
\]
The six \(K_7\) triples are
\[
 (5,6,1),(6,5,1),(9,20,13),(11,24,13),(10,17,1),(12,20,8),
\]
and the four \(K_9\) triples are
\[
                     (6,9,3),(9,18,9),(9,6,1),(12,27,17).
\]
They follow by multiplication in the displayed full defining bases.

The anchors have odd valuations respectively at the primes
\(127,191,631,19\):
\[
 1016=8(127),\quad764=4(191),\quad5048=8(631),\quad
 16872=8\cdot3\cdot19\cdot37.
\]
NONE of the other time norms in its field row is divisible by
that prime.  The exact nonzero residues are respectively
\[
 (53,74,40,45,36),\quad(185,93,13,108),\quad
 (67,234,127,522,44),\quad(10,11,17).
\]
Thus every UNEQUAL anchor/model time ratio has nonsquare rational
norm and cannot be a square in \(K\).
No square norm is used as a sufficient test, and no native,
orthogonal, parallel or dotted mutual case is omitted by this
necessary invariant.

Only \(V=U\) at the actual finite anchor remains.
The endpoint walls are distinct, and their canonical raw norms
are \(2U\).
Equal time squares give equal POSITIVE physical times;
their spatial product is \(U^2/17\).
Therefore the actual two-cycle is
\[
 g={U(34-15U)\over17},\qquad
 L={g^2\over U^2}=\left({34-15U\over17}\right)^2
                     \in\mathcal O_K .
\]
The zero branch would force the irrational anchor \(U=34/15\),
impossible, but integral closedness applies in either case.
It makes \(15U/17\), hence \(U/17\), integral:
\(15\cdot8-17\cdot7=1\).
The four anchor norms are \(2,1,1,3\), in degrees \(2,2,3,3\);
their quotients by the corresponding powers of seventeen
are not integers.
This final contradiction excludes every diagonal pair too.

The whole-\(A_{16}\) exclusion is complete: actual endpoint
lengths and helper/golden cases, BOTH isolated, degree at most six,
full two-dotted interval, finite anchors and actual fields, complete
full-ring domains, every unequal time ratio, and the actual
diagonal two-cycle.
It is manual at the printed finite inequalities and norm tables,
with no new enumeration-output premise, altered dimension bound
or assumed equality of disconnected root scales.

\section{One actual rank-one opposite}
\label{app:arithrankone}

For the full forests \(A_{15}+A_1,B_{15}+A_1,D_{15}+A_1\),
the correct first split concerns the SAME actual opposite \(d\)
of the old rank-one direction \(e\).
Either it contacts the old large component, or all those products
vanish.  In the isolated case physical unit normalization gives
\[
                 d=\sqrt h\,p-\sqrt{1+h}\,e,\qquad h>0 .
\]
No unrelated wall is chosen in the second case, and this identity
does not identify disconnected field-root norm squareclasses.
We first close the whole \(A_{15}+A_1\) case and record its
separate \(B_{15}\) nonisolated corollary.

\subsection{Nonisolated rank-one opposites beside \(A_{15}\)}
\label{app:arithA15nonisolated}

Let the original full forest be \(A_{15}+A_1\), with old
unit \(A\)-roots in zero-sum \(\mathbb R^{16}\), and write
\[
 d=Ap-v_A-\gamma e,\quad A>0,\quad\gamma>0,\quad
 Q=\|v_A\|^2,\quad A^2=\gamma^2+Q-1.
\]
At the actual endpoint of the rank-one edge, old \(A_{15}+d\)
is a finite rank-sixteen white diagram.
Its COMPLETE nonisolated extensions are
\[
\begin{array}{c|c|c}
\text{role}&v_A&Q\\ \hline
\text{endpoint three}&\omega_j/2,\ j=1,15&15/32\\
\text{endpoint four}&\omega_j/\sqrt2,\ j=1,15&15/16\\
\text{next-to-end three}&\omega_j/2,\ j=2,14&7/8 .
\end{array}
\]
These are \(A_{16},B_{16},D_{16}\), including both orientations.
Two old-path contacts make a nonpositive cycle; no exceptional
rank-sixteen completion contains this long path.

We must have \(\gamma>1\).
If \(\gamma<1\), the full positive actual replacement
\[
                 r=p+{A\over1-\gamma^2}(d+\gamma e)
\]
has negative norm, positive time, zero products with \(e,d\)
and all retained old \(A\)-roots except the contacted node \(j\).
Every unknown wall is nonpositive termwise, and its omitted
\(n_j\) product is strictly negative.
It is therefore the ACTUAL endpoint of old edge \(j\).
The SAME \(d\) would close both that direction and the
orthogonal old rank-one direction, contradicting the
original-farthest orthogonal obstruction.
For \(\gamma=1\), \(d+e\) is an actual future null.
Thus \(Q<A^2\) and the entire surviving own domain is dotted.

The actual rank-one edge length is
\(C=\gamma/\sqrt{1-Q}\), with energy at least \((C-1)/(C+1)\).
At \(Q\ge7/8\), \(C>\sqrt8>7/3\) gives energy \(>2/5\).
Together with the independently proved same-vector
\(e_{A_{15}}>.8348\), this excludes endpoint-four and
next-three throughout their unbounded own domains.
For endpoint-three, if \(\gamma\ge41/40\),
\[
                  C^2\ge1681/850>49/25
\]
gives rank-one energy \(>1/6\), and
\(.8348+1/6=1+11/7500\).
Only \(1<\gamma<41/40\) remains.

Here the ACTUAL positive core
\[
                  d+\gamma e=Ap-\omega_j/2
\]
is future timelike of square \(1-\gamma^2\), zero on \(e\),
negative on \(d\), and nonpositive on every other original wall.
Original maximality gives
\(h_j/2\ge A-\sqrt{\gamma^2-1}\).
The right side decreases with \(\gamma\).
At \(41/40\), \(A^2=831/1600>.72^2\) and
\(\sqrt{\gamma^2-1}=9/40\).
Thus the appropriate ACTUAL signed \(A_{15}\) endpoint satisfies
\[
                         x_{\rm end}>\sqrt2(.72-.225)>.69 .
\]

Rebuild its reference before using the new pressure.
Take the common N16 raw heights
\(2.021,2.137,2.205,2.227\), reflected,
with their possible single golden hole.
Internal omission loses at most \(.000576\).
End omission replaces the first tail by \(0\)--\(6\), giving
\[
                         E_{\rm end}=48050651/60000000>.80084 .
\]
Thus every RAW reference has energy \(>.79\), before any
earlier golden-endpoint graft.
At the selected endpoint its flat block has length five or six,
mean \(.25<b<.3\), and next coordinate \(<.101\), including
an adjacent internal omission.
Graft this block afresh to \(.69\).
The other coordinates are \(>(L(.25)-.69)/(L-1)\ge.14\);
the changed central coefficient is
\(>(.14-.101)/\sqrt2>.027>.022\).
Every other positive central coefficient retains \(.022\);
the golden hole remains coefficient zero.
The exact sum-preserving zero-sum change has zero old block cross,
so the new ONE reference has energy
\[
                    >.79+(6/5)(.69-.3)^2=.97252 .
\]
Absorb EVERY retained actual isolated foot again into THIS
reference, using the common N16 loss inequality and its NEW
coefficients.
This gives \(e_A>.97252\), including adjacent and foreign feet.
No earlier optimized endpoint gain or foot energy was added.

Finally the original rank-one edge has
\(C>\sqrt{32/17}>11/9\), so its distinct orthogonal energy is
\(>1/10\).
Adding that energy ONCE gives
\(\|x\|^2>.97252+.1>1\).
This excludes the entire nonisolated actual rank-one opposite
beside \(A_{15}\), over every actual field.

For later use the corresponding \(B_{15}\) argument is separate.
A finite rank-sixteen extension of old \(B_{15}\) can only be
\(B_{16}\), attaching by three at its LONG endpoint.
Thus \(v_B=\Omega_1/2\), \(Q=1/2\), and
\(A^2=\gamma^2-1/2\).
The same reciprocal actual projection proves \(\gamma>1\).
The original rank-one energy is then
\(>(\sqrt2-1)/(\sqrt2+1)=3-2\sqrt2>.17\).
The already written fresh \(B_{15}\) bulk in the proof of
Proposition~\ref{prop:f16}, with arbitrary spectators and both
initial-five cases, gives \(e_B>E_{13}>.88989\).
Their sum exceeds \(1.05989\).
Thus a nonisolated actual rank-one opposite is impossible
beside \(B_{15}\) as well.  No \(D_{15}\) conclusion is
transported from this corollary.

\subsection{The isolated \(A_{15}\) case and its invariant cycles}
\label{app:arithA15isolatedcycles}

Now suppose the SAME actual opposite is isolated:
\(d=\sqrt h\,p-\sqrt{1+h}\,e\).
If \(h\ge1\), its actual edge has energy at least
\(3-2\sqrt2>.17\); together with \(e_A>.8348\) this
already contradicts the norm budget.
Assume \(0<h<1\) and put \(V=4h\).
The actual two-cycle makes \(V\) integral and all nonphysical
places satisfy \(-4<V^\sigma<0\).
Physically \(0<V<4\).
Consequently \(V/4\) CANNOT be integral, since its nonzero norm
would have magnitude below one.
The actual field has degree at least two by the stated
Meyer/compact-anisotropy input for the ENTIRE rank-seventeen form.

We need pressures only at
\[
 J=\{5,6,7,9,10,11\},\qquad
 \ell_i=.49\ (i\text{ odd}),\quad
 \ell_6=\ell_{10}=.496 .
\]
The midpoint eight will never be assigned a pressure or coefficient.

Let an actual central opposite be
\[
 c=Ap-v_A-\beta e,\quad Q=\beta^2<1,\quad
 \lambda=-c\cdot d=A\sqrt h-\beta\sqrt{1+h}\ge0.
\]
The INVARIANT actual two/three cycles are
\[
 L=4\lambda^2,\quad W=4\beta^2,\quad
T=8\beta\sqrt{1+h}\lambda,\qquad L,W,T\in\mathcal O_K .
\]
Zero \(\beta\) or \(\lambda\) is included.
Direct expansion cancels the mixed radical:
\[
                  L+T=VA^2-W(1+V/4).
\tag{Rankonecycle}\label{eq:arithrankonecycle}
\]
No individual unit coefficient or disconnected raw time is
assumed \(K\)-valued.

For an isolated-within-\(A\) cap,
\(v_A=\gamma\omega_i\), \(B=4\gamma^2\),
this gives the exact full-foreign identity
\[
              L+T+V+W={i(16-i)BV\over32}\in\mathcal O_K.
\tag{A15isocycle}\label{eq:arithA15isocycle}
\]
At EVERY nonphysical embedding, positive definiteness on the
whole old \(A_{15}\) and this cap gives
\(0<B^\sigma<4/D_i\); its additional foreign square is
nonnegative and only strengthens the bound.
Together with \(-4<V^\sigma<0\), this gives
\(\lvert B^\sigma V^\sigma\rvert<16/D_i\).
For one white-three contact,
\(v_A=\gamma\omega_i+\omega_j/2\),
\(t=2\gamma\), and
\[
 M=i(16-i)t^2+2\min(i,j)(16-\max(i,j))t+j(16-j)-32,
\]
the corresponding identity is
\[
                         L+T+W=MV/32\in\mathcal O_K.
\tag{A15whitecycle}\label{eq:arithA15whitecycle}
\]
Thus the foreign radical is retained through its triangle term.
Dropping that term would not prove these congruences.

\subsection{Finite native parameters and all dyadic places}
\label{app:arithA15dyadic}

For \(a=2\cos(\pi/m)\), \(m\ge4\), in any containing field,
every dyadic valuation normalized by \(v(2)=1\) satisfies
\[
                         0\le v(a),v(a+1)\le1/2 .
\]
Here is the full root-of-unity argument.
Write \(a=\zeta+\zeta^{-1}\), \(\zeta\) primitive \(2m\)-th.
Prime-to-two roots keep their order on reduction: a nontrivial
odd-order root reducing to one would make its zero geometric
sum reduce to a nonzero odd integer.
Thus \(v(a)>0\) occurs only when \(2m=2^s\).
For \(s\ge3\),
\(v(a)=v(\zeta^2+1)=1/2^{s-2}\le1/2\).
The formula follows because conjugate primitive \(2^r\)-root
differences have equal valuation and product two.
Likewise
\[
                a+1=\zeta^{-1}(\zeta^3-1)/(\zeta-1)
\]
has positive valuation only when \(2m=3\cdot2^s\);
for \(m\ge4\), \(s\ge2\), and
\(v(a+1)=1/2^{s-1}\le1/2\).
For \(m=3\), \(a=1\) and \(v(a+1)=1\); that exception is
removed by an actual affine cycle BEFORE this lemma is used.
In particular \(B=4\cos^2(\pi/m)=a^2\) satisfies \(v(B)\le1\)
for every finite \(m\ge3\).
If \(a\notin K\), use the cyclotomic compositum and descend
the normalized valuation of \(B\in K\).
These bounds persist under every actual field extension.

For a finite one-white-three own, the actual three-cycle
along the old path excludes own three.
For every other finite own, apply the bounds to \(t\).
Put \(a_i=i(16-i)\) and \(C=\min(i,j)(16-\max(i,j))\).
The following four parity cases prove \(v(M)\le3\)
at EVERY dyadic place.

If \(i,j\) are odd, rewrite in \(s=t+1\):
\[
 M=a_is^2+2(C-a_i)s+[r(16-r)-32],\qquad r=|i-j|.
\]
The last two coefficients are divisible by four and \(a_i\)
is odd.  The leading term has valuation at most one, while
the others have valuation at least two.
If \(i\) is odd and \(j\) even, use \(t\) instead:
\(a_i\) is odd, \(C\) even and \(a_j\) divisible by four.
The same unique-minimum argument gives \(v(M)\le1\).
If \(i\) is even and \(j\) odd, the constant is odd and the
other coefficients are even, giving \(v(M)=0\).
If both are even, \(i=6,10\) makes \(k=i/2\) odd.
With \(l=j/2\),
\[
 M/4=k(8-k)t^2+2\min(k,l)(8-\max(k,l))t+l(8-l)-8 .
\]
For even \(l\) use the \(t\)-argument; for odd \(l\) use
\(s=t+1\), whose constant
\(r(8-r)-8\), \(r=|k-l|\), is divisible by four.
Thus \(v(M/4)\le1\) and \(v(M)\le3\).
No midpoint-eight assertion has been made.

Equation~\eqref{eq:arithA15whitecycle} now forces \(v(V)\ge2\)
at every dyadic place.
At all other finite places \(V\) is already integral.
Hence \(V/4\) would be integral, impossible.
This excludes EVERY finite one-white-three cap at the six
required sites, including foreign contacts and both orientations.
The same proof applies to a dotted \(t\) whenever both
\(v(t),v(t+1)\le1/2\) have been established.

\subsection{The fresh six-site inventory at losses \(.49,.496\)}
\label{app:arithA15freshinventory}

At each \(i\in J\), both retained old paths have ranks at
least four.
Their COMPLETE finite-white list is no old-\(A\) contact;
one three at an admissible end or \(D/E\) position;
one four at a path end; or one endpoint three in each path.
The two-path bridge has rank fifteen, so it is \(A_{15}\):
a third arm through the new root would have two old arms at
least four, incompatible with finite \(D\), and no exceptional
finite type has this rank.
The old \(e\)-contact \(\beta\) is retained throughout.
It only restricts the white list further.

Own zero is not silently discarded.
Without old-\(A\) contact its time \(Q-1\) is negative.
For a white three at an internal old node, the full old diagram
and cap have actual affine arms \((1,2,5)\) or \((1,3,3)\),
unless the contact is at \(1,2,14,15\).
At global ends one or fifteen, the full finite white path
has rank at least five; adding \(e\) allows only
\(\beta=0,1/2,1/\sqrt2\).
Then \(A^2=15/32+Q-1<0\).
At two or fourteen, \(\beta=0\) gives time \(-1/8\).
A positive foreign contact in that finite branched white
component forces \(\beta=1/2\), and the full actual roots
then have the affine arms \((1,2,5)\).
An internal white four has the \(4/3/3\) star.
At a global end white four, any positive \(e\)-contact would
put a three BEFORE the four in a path of rank at least six,
or make a nonfinite \(4,4\) principal; hence \(Q=0\) and
time \(15/16-1<0\).
Two white threes at own zero give the actual all-three cycle.
Every null is a positive sum of actual roots, future through
the cap time, and nonpositive against ALL omitted walls.

For every positive own, the parallel value is excluded BEFORE
the remaining case dispatch: \(c+n_i\) is its actual feasible
future null, including positive foreign contacts.
For white three on a side of order \(m=i\) or \(16-i\),
number its node \(j\) from that side's far end.
The actual residual and numerator are
\[
 \delta=1-j(m-j)/(2m),\qquad \sigma=\gamma+j/(2m).
\]
For white four at an end they are
\(\delta=1/m\), \(\sigma=\gamma+j/(\sqrt2m)\).
For two white threes, \(\delta=1/D_i\),
\(\sigma\ge\gamma+1/D_i\).
The ENTIRE old-\(e\) cost changes the denominator to
\(\delta-Q\), and
\[
                   C_i={\sigma\sqrt{D_i}\over\sqrt{\delta-Q}} .
\]
If the desired \(\ell_i\)-pressure fails, this implies
\[
 {\sigma\over\sqrt{\delta-Q}}
       \le R_{\ell_i}(\sqrt{D_i})<1.881 .
\tag{A15ratio}\label{eq:arithA15ratio}
\]
For completeness the ratio target increases with \(q\) and
decreases with \(\ell\) on \(q\le\sqrt8<2.829\),
\(\ell\ge.49\).
At \(q=2.829,\ell=.49\), positive denominator clearing gives
the strict margin119649968.
Thus this is a fresh changed-loss inventory, rather than
reuse of the old \(.601\) weak list.

Finite one-white-three owns have already been excluded by
the invariant parity argument, including own three's affine cycle.
Thus only DOTTED \(\gamma>1\) remains there.
At a \(D\)-position the ratio is
\(>(1+1/m)/\sqrt{2/m}\ge6/\sqrt{10}>1.881\).
At an \(E\)-position \(\delta\le1/4,\sigma>1\), stronger still.
Foreign cost decreases the denominator and cannot weaken this.
For near white four, every positive own gives ratio
\[
            \ge(1/2+4/(5\sqrt2))\sqrt5>2.38 .
\]
For far white four, owns three, four and six have actual
affine principals using old roots at most three steps on each
side of \(i\), avoiding the helper at one or fifteen.
The remaining own is at least \(\varphi/2>.809\), giving
\(.809\sqrt5+1/\sqrt{10}>2.12\).
For two white threes, own three has an actual cycle.
Every other positive own is at least \(1/\sqrt2>.7\),
and its ratio exceeds \(.7(2.6)+.35>1.881\).
Every finite-white role except dotted near/far three has now
been included and removed.

\subsection{Every remaining dotted helper fails over every field}
\label{app:arithA15dottedhelpers}

For near white three the ratio is
\[
 {\,\gamma+(m-1)/(2m)\over\sqrt{(m+1)/(2m)-Q}}.
\]
At \(m\ge7,\gamma>1\), it exceeds \(10/\sqrt{28}>1.881\).
For \(m=5,6\), weakness gives
\(2<t=2\gamma<2.11555<53/25\).
These are \(i=5,j=4\), \(i=6,j=5\), and reflections.
Lift EVERY nonidentity embedding of \(\mathbb Q(t)\) to \(K\).
Strict full-old-\(A\) positivity gives respectively
\[
 (55/8)t_\sigma^2+11t_\sigma+2<0,\qquad
 (15/2)t_\sigma^2+(25/2)t_\sigma+23/8<0 .
\]
The foreign square is nonnegative, strengthening these conditions.
Both negative quadratic intervals lie in \((-1.4,-.2)\);
check positive endpoint values at \(-7/5,-1/5\).

If the minimal degree is at least three and some small conjugate
is at least \(-1\), the nonzero integral norm of \(t(t+1)\)
has magnitude
\[
              <6.6144(1/4)(.56)^{d-2}\le.926016<1 .
\]
If all small conjugates are below \(-1\), the norm of \(t+1\)
has magnitude \(<3.12(.4)^{d-1}<1\).
In degree two the integral trace is one, and the norm
\(t(1-t)\) lies strictly between \(-3,-2\).
In degree one no integer lies in \((2,2.12)\).
Thus every weak near dotted helper is impossible, including
foreign contacts and proper actual fields.

For far white three, if \(\gamma\ge34/25\) its squared ratio
is at least its \(m=5\) value \(5329/1500>1.881^2\).
To check the minimum, put \(r=1/m\).
The derivative of \(2(\gamma+r/2)^2/(1+r)\) has sign
\(1-\gamma+r/2<0\), and the ratio increases with \(\gamma\).
Hence a weak far helper has \(2<t<68/25\).
Conjugate positivity gives the entire strip
\[
                            -4/5<t_\sigma<12/25 .
\]
The N16 endpoint cost checks at \(-.8,.48\) have margins
\(3/40,59/1000\); they increase over the reflected half-window.
Again positive foreign cost only strengthens them.

For \(f(t)=t(t^2-1)\), small conjugate magnitude is at most
\(2/(3\sqrt3)\).
Degree at least four would give norm magnitude
\[
 {68\over25}\left\{\left({68\over25}\right)^2-1\right\}
                         {8\over81\sqrt3}<1 .
\]
For an explicit positive seal, eight times the physical factor
is139.229184, below \(81(1.732)=140.292\).
Thus the minimal degree is at most three.
At degree three, \(|\mathrm N(t)|<2.72(.8)^2<2\), so it is
a unit.
Two positive small conjugates would give norm below one.
For opposite signs their absolute product \(1/t\) forces
\[
 t>125/48,\quad
 t_{\rm positive}>125/272,\quad
 |t_{\rm negative}|>625/816,
\]
putting the integral trace strictly between two and three.
For two negative conjugates, the trace is one and \(t<2.6\).
The integral second coefficient
\(1/t-t(t-1)\) lies in \((-4,-1.5)\), hence is \(-2\) or \(-3\).
The two cubics are
\(X^3-X^2-2X-1\), of discriminant \(-31\), and
\(X^3-X^2-3X-1=(X+1)(X^2-2X-1)\).
Neither is the minimal polynomial of the totally real cubic.
Degree three is excluded.

Degree two has trace two or three.
Trace two has norm in \((-1.9584,0)\), hence \(-1\);
trace three has positive norm below two, hence one.
The exact possibilities are
\[
                         t=1+\sqrt2,\qquad t=\varphi^2 .
\]
Degree one is the forbidden parallel two.
For \(1+\sqrt2\), \(t\) is a dyadic unit and
\(v(t+1)=1/2\).
For \(\varphi^2\), both \(t\) and \(t+1=\varphi+2\) have
odd norm and are dyadic units.
These facts persist in every actual field extension.
Thus the already proved four-parity argument applies to \(M\)
for BOTH remaining dotted possibilities.
Equation~\eqref{eq:arithA15whitecycle} forces \(V/4\) integral,
a contradiction.
All weak dotted helpers are excluded.

\subsection{The isolated central caps through the full \(h<1\) interval}
\label{app:arithA15centralisolated}

For isolated-within-\(A\) contact, use
\eqref{eq:arithA15isocycle}.
At odd sites \(i(16-i)\) is coprime to32, so integer Bezout
gives \(BV/32\in\mathcal O_K\).
At six and ten, the gcd is four, giving \(BV/8\in\mathcal O_K\).
For ANY finite own the dyadic bound \(v(B)\le1\) then forces
\(v(V)\ge4\) or \(\ge2\), respectively.
Thus every finite isolated own would make \(V/4\) integral,
impossible; own zero has negative time.
Parallel was already excluded globally by its actual null.

For dotted own, \(u=(1-Q)/\gamma^2\in(0,1)\).
Its positive ACTUAL core
\(c+\beta e+\gamma u n_i\)
has square \(-\gamma^2u(1-u)\), positive time, zero \(c,e\)
products, negative \(n_i\) product and all other wall products
nonpositive.
It gives
\[
 h_i>\sqrt{D_i-u}-\sqrt{u(1-u)},\qquad
 \text{loss}\le au+\sqrt{u(1-u)} ,
\]
with \(a<.2\) at the odd sites and \(a<.19\) at six and ten.
If \(u\le1/4\) at an odd site, the loss is
\(<1/20+\sqrt3/4<39/80<.49\).
Thus an odd weak cap forces \(u>1/4\) and \(B<16\).
The integral \(BV/32\) has physical magnitude \(<2\) and
EVERY other magnitude \(<1/(2D_i)\le4/55\).
Since the actual degree is at least two, its nonzero norm has
magnitude \(<8/55<1\), impossible.

At six or ten, for \(u\le.27\) the loss is
\(<.0513+.444=.4953<.496\).
An even weak cap therefore has \(u>.27\), \(B<400/27\).
The integral \(BV/8\) has physical magnitude \(<200/27\),
and all other magnitudes \(<4/15\).
At degree at least three its norm is
\(<(200/27)(4/15)^2=128/243<1\).
At degree two, instead use the rational integer norms:
\[
       0<\mathrm N(B)<640/81<8,\qquad
       0<|\mathrm N(V)|<16 .
\]
The first is positive and the second negative before taking
absolute values.
Integrality of \(BV/8\) requires64 to divide their product.
But an integer from one through seven has dyadic valuation
at most two, and one through fifteen at most three.
Their sum is less than six, a contradiction.
This argument permits arbitrary ramification and full maximal order.

Combining the complete own-zero, finite-white, finite, parallel,
dotted-helper and isolated cases proves, for EVERY \(i\in J\),
\[
                            h_i>\sqrt{D_i}-\ell_i .
\]
There is still NO midpoint-eight assertion.

\subsection{One final reference for every isolated rank-one height}
\label{app:arithA15isolatedcomposition}

Choose the rational heights
\[
 H_5=H_{11}=2.132,\quad H_6=H_{10}=2.2426,\quad
 H_7=H_8=H_9=2.3162
\]
and linear outer tails.
Their six positive constraints lie strictly below the just
proved pressures by direct squares at \(55/8,15/2,63/8\).
The midpoint coefficient is zero, the slopes decrease, and
\[
                 E=2.132^2/5+.1106^2+.0736^2=.92673412 .
\]
For \(1/2\le h<1\), the separate rank-one edge energy is
at least \(5-2\sqrt6>.1\); hence the original norm exceeds one.
Together with the earlier \(h\ge1\) argument this leaves
only \(0<h<1/2\).

We retain BOTH actual \(A_{15}\) endpoint pressures in that case.
Put
\[
 T=\sqrt{1+4h}-2\sqrt h,\qquad
 u={\sqrt h\over\sqrt{1+h}+1},\qquad
 \alpha=\sqrt{15/16}.
\]
The actual rank-one edge gives energy \(\ge u^2\), \(0<u<1/3\),
and direct substitution gives
\((1-u^2)(1-T^2)=8uT\).
The rank-free actual-\(I\) wall comparison
\eqref{eq:Ibarrier} applies with every foreign contact retained.
Explicitly, for another actual opposite with time \(A\)
and old-\(e\) contact \(\beta=0\) or \(1/2\le\beta<1\),
its chosen spatial derivative square is at most
\(1+A^2-\beta^2\).
For \(\beta=0\), the native gap gives \(A^2\ge1/(4h)\);
for \(\beta>0\), acuteness gives
\(A^2\ge(1+h)\beta^2/h\).
Both imply full tangent square \(\ge1/(1+4h)\).
In the positive-contact case the comparison, after positive
denominators are cleared, has numerator
\(h((5+4h)\beta^2-1)>0\).
Thus its normalized half-edge height is at least \(T\).
At BOTH signed \(A_{15}\) endpoints this gives
\(x_1,-x_{16}\ge\alpha T\).
This is the original wall comparison, not a new field assumption.

The first and last five-coordinate reference means are
\(b=2.132/(5\sqrt2)<.30152<.302\).
Since \(T>\sqrt3-\sqrt2>.317\) and \(\alpha>.968\),
\(\eta=\alpha T>b\).
Graft BOTH disjoint blocks to \(\eta\).
Their other coordinates
\((5b-\eta)/4>1/8\) exceed the next old coordinate
\(.1106/\sqrt2<.079\).
Thus all dual coefficients remain nonnegative, supported only
on the six proved pressures and the two ACTUAL endpoint bounds;
the midpoint coefficient remains zero.
The single rebuilt reference gives
\[
                         e_A>E+(5/2)(\eta-b)^2 .
\]

For completeness the scalar barrier \(T+3u/2>.76\) follows
directly from the displayed relation.
At \(T_0=.76-1.5u>0\), its left-minus-right is bounded below by
\[
                  264/625-(19/5)u+(21419/2500)u^2>0 .
\]
The quadratic discriminant is \(-55964/1562500<0\);
the cubic term was bounded using \(u<1/3\).
The relation decreases strictly with positive \(T\), proving
the barrier.
With \(D=\eta-b>0\),
\[
 D+(3/2)\alpha u>.76\alpha-b>.43368.
\]
Weighted Cauchy gives
\[
 (5/2)D^2+u^2>
 {.43368^2\over2/5+(9/4)\alpha^2}
 ={.43368^2\over803/320}>.074 .
\]
The positive comparison is
\(.1880783424>.18569375\).
Adding the original rank-one energy ONCE to the rebuilt
\(A_{15}\) reference yields
\[
                     \|x\|^2>.92673412+.074=1.00073412>1 .
\]
All \(h>0\), including boundaries \(1/2,1\), are excluded
for the isolated opposite.

Together with the complete nonisolated argument, this excludes
the ENTIRE original \(A_{15}+A_1\) family, with either ordinary
nonshort rank one or actual short \(I\).
The dichotomy concerns one actual wall throughout, and the two
case contradictions were never added to one another.
No historical frontier-count subtraction is needed.
The separate \(B_{15}\) nonisolated corollary only supplies its
isolation reduction; its isolated ordinary and actual-\(I\)
branches remain to be proved.









